\documentclass[11pt,oneside]{article}

% =====================================================================
%  Asymptotic transseries for A000081 and its inverse
%  Self-contained research article prepared for the ProveIt project.
% =====================================================================
\usepackage[a4paper,margin=25mm,headheight=15pt]{geometry}
\usepackage[T1]{fontenc}
\usepackage[utf8]{inputenc}
\IfFileExists{libertinus.sty}{\usepackage{libertinus}}{\usepackage{lmodern}}
\usepackage{microtype}
\usepackage{amsmath,amssymb,amsthm,mathtools,mathrsfs,bm}
\usepackage{booktabs,longtable,tabularx,array,multirow}
\usepackage{graphicx}
\usepackage{xcolor}
\usepackage{enumitem}
\usepackage{fancyhdr}
\usepackage{titlesec}
\usepackage{float}
\usepackage{xurl}
\usepackage{listings}
\usepackage[most]{tcolorbox}
\usepackage{hyperref}
\usepackage[nameinlink,noabbrev,capitalise]{cleveref}

\definecolor{linkblue}{RGB}{24,72,124}
\definecolor{deepgreen}{RGB}{23,102,69}
\definecolor{deepred}{RGB}{140,42,42}
\definecolor{deepgold}{RGB}{122,91,19}
\definecolor{deepviolet}{RGB}{89,55,128}
\definecolor{shadegray}{RGB}{246,247,249}
\definecolor{rulegray}{RGB}{195,201,208}
\definecolor{palered}{RGB}{251,236,236}
\definecolor{palegold}{RGB}{251,245,230}
\definecolor{palegreen}{RGB}{233,245,243}
\definecolor{paleblue}{RGB}{237,244,251}
\definecolor{paleviolet}{RGB}{244,239,250}

\hypersetup{
  hypertexnames=false,
  colorlinks=true,
  linkcolor=linkblue,
  citecolor=deepgreen,
  urlcolor=linkblue,
  pdftitle={Asymptotic Transseries of the Butcher--Polya Rooted-Tree Numbers and Their Inverse},
  pdfauthor={Research article prepared for Vladimir Reshetnikov and the ProveIt project},
  pdfsubject={Full asymptotic expansion of OEIS A000081, Bell-polynomial coefficient formulae, Lambert-W functional inversion, and exponentially small singularity sectors},
  pdfkeywords={Butcher trees, Polya trees, rooted trees, A000081, Otter constant, asymptotic expansion, transseries, inverse asymptotics, Lambert W, Bell polynomials, Lagrange inversion}
}

\pagestyle{fancy}
\fancyhf{}
\fancyhead[L]{\small Butcher--P\'olya tree transseries}
\fancyhead[R]{\small\nouppercase{\leftmark}}
\fancyfoot[C]{\thepage}
\renewcommand{\headrulewidth}{0.4pt}

\setcounter{tocdepth}{2}
\setcounter{secnumdepth}{3}
\setlength{\parindent}{0pt}
\setlength{\parskip}{0.46em}
\setlist[itemize]{leftmargin=1.45em,itemsep=0.25ex,topsep=0.55ex}
\setlist[enumerate]{leftmargin=1.85em,itemsep=0.35ex,topsep=0.55ex}
\allowdisplaybreaks[2]
\emergencystretch=2em
\numberwithin{equation}{section}

\theoremstyle{plain}
\newtheorem{theorem}{Theorem}[section]
\newtheorem{proposition}[theorem]{Proposition}
\newtheorem{lemma}[theorem]{Lemma}
\newtheorem{corollary}[theorem]{Corollary}
\newtheorem{conjecture}[theorem]{Conjecture}
\theoremstyle{definition}
\newtheorem{definition}[theorem]{Definition}
\newtheorem{example}[theorem]{Example}
\newtheorem{algorithm}[theorem]{Algorithm}
\theoremstyle{remark}
\newtheorem{remark}[theorem]{Remark}
\newtheorem{warning}[theorem]{Warning}

\newtcolorbox{mainresultbox}[1][]{%
  enhanced,breakable,colback=paleblue,colframe=linkblue,
  boxrule=0.75pt,arc=1.2mm,left=2.2mm,right=2.2mm,top=1.5mm,bottom=1.5mm,
  title={Main result},fonttitle=\bfseries,#1}
\newtcolorbox{summarybox}[1][]{%
  enhanced,breakable,colback=palegold,colframe=deepgold,
  boxrule=0.6pt,arc=1.2mm,left=1.8mm,right=1.8mm,top=1.2mm,bottom=1.2mm,
  title={Section summary},fonttitle=\bfseries,#1}
\newtcolorbox{checkpointbox}[1][]{%
  enhanced,breakable,colback=palegreen,colframe=deepgreen,
  boxrule=0.6pt,arc=1.2mm,left=1.8mm,right=1.8mm,top=1.2mm,bottom=1.2mm,
  title={Structural checkpoint},fonttitle=\bfseries,#1}
\newtcolorbox{warningbox}[1][]{%
  enhanced,breakable,colback=palered,colframe=deepred,
  boxrule=0.6pt,arc=1.2mm,left=1.8mm,right=1.8mm,top=1.2mm,bottom=1.2mm,
  title={A point requiring care},fonttitle=\bfseries,#1}
\newtcolorbox{frontierbox}[1][]{%
  enhanced,breakable,colback=paleviolet,colframe=deepviolet,
  boxrule=0.65pt,arc=1.2mm,left=1.8mm,right=1.8mm,top=1.2mm,bottom=1.2mm,
  title={Beyond-all-orders completion},fonttitle=\bfseries,#1}

\lstdefinestyle{pseudo}{%
  basicstyle=\ttfamily\small,columns=fullflexible,frame=single,
  rulecolor=\color{rulegray},backgroundcolor=\color{shadegray},
  showstringspaces=false,keepspaces=true,breaklines=true,
  xleftmargin=0.4em,xrightmargin=0.4em,aboveskip=0.8em,belowskip=0.8em}
\lstset{style=pseudo}

% Document-local notation.
\newcommand{\e}{\mathrm e}
\newcommand{\ii}{\mathrm i}
\newcommand{\dd}{\,\mathrm d}
\newcommand{\Ord}{\mathcal O}
\newcommand{\littleo}{\mathrm o}
\newcommand{\N}{\mathbb N}
\newcommand{\Nzero}{\mathbb N_0}
\newcommand{\C}{\mathbb C}
\newcommand{\R}{\mathbb R}
\newcommand{\Tcal}{\mathcal T}
\newcommand{\Cay}{\mathcal C}
\newcommand{\Bhat}{\widehat{B}}
\newcommand{\Coeff}[2]{[\,#1^{#2}\,]}
\newcommand{\fall}[2]{(#1)_{#2}}
\newcommand{\Log}{\operatorname{Log}}
\newcommand{\Li}{\operatorname{Li}}
\newcommand{\Res}{\operatorname*{Res}}
\newcommand{\sgn}{\operatorname{sgn}}
\DeclareMathOperator{\supp}{supp}
\DeclareMathOperator{\val}{val}

\title{\Huge\bfseries Asymptotic Transseries of the\\[1mm]
Butcher--P\'olya Rooted-Tree Numbers\\[2mm]
\Large Bell-Polynomial Coefficients, Lambert--$W$ Reversion,\\
and Exponentially Small Singularity Sectors}
\author{Research article prepared for Vladimir Reshetnikov\\and the ProveIt project}
\date{September 2026}

\begin{document}
\maketitle

\begin{abstract}
Let $a_n$ denote OEIS A000081, the number of unlabelled rooted non-plane
 trees with $n$ vertices, also called P\'olya trees and, in the numerical-
analysis literature, the unlabelled rooted trees underlying Butcher series.
Its ordinary generating function $T(z)=\sum_{n\ge1}a_nz^n$ satisfies
\[
 T(z)=z\exp\!\left(\sum_{k\ge1}\frac{T(z^k)}{k}\right).
\]
This article develops an all-orders asymptotic and transseries calculus for
$a_n$ and for its large-value functional inverse.  The dominant singularity
$\rho$ is put into the Cayley normal form $T=\Cay\circ\zeta$, and every
Puiseux coefficient at $z=\rho$ is expressed by a finite ordinary-Bell-
polynomial formula in the derivatives of $\zeta$.  Exact coefficient
transfer is then expanded by Bernoulli polynomials and complete exponential
Bell polynomials.  The resulting compact master formula is
\[
 a_n\sim \frac{\rho^{-n}}{\sqrt{\pi}\,n^{3/2}}
 \sum_{m\ge0}\frac{\tau_m}{n^m},\qquad
 \tau_m=\sum_{j=0}^{m}
 \frac{(-1)^{j+1}(2j+1)!!}{2^{j+1}}
 t_{2j+1}\,g_{m-j}\!\left(j+\tfrac12\right),
\]
where the $t_r$ are the local Puiseux coefficients and the $g_r(\beta)$
are given explicitly by complete Bell polynomials in Bernoulli-polynomial
increments.  Thus all asymptotic coefficients are obtained by a finite,
triangular, auditable calculus.

Writing the dominant model as
$C\e^{\lambda x}x^{-3/2}S(x^{-1})$, its exact carrier inverse is a real
$W_{-1}$ expression.  If $H=\log S$, the complete algebraic correction is
$\nu(y)=x_0(y)+\sum_{m\ge1}d_mx_0(y)^{-m}$, where the $d_m$ obey an explicit
triangular Bell-polynomial recurrence.  An equivalent polynomial--logarithmic
form $\nu(y)=X+q\log X+\sum_{m\ge1}P_m(\log X)X^{-m}$ is derived, again with
a general coefficient formula.  Finally, a singular-germ recursion produces
the beyond-all-orders sectors associated with inherited root singularities
and any additional critical preimages; a corresponding carrier-first
implicit recursion gives the inverse transseries.  We distinguish rigorously
between the smooth functional inverse, the stepwise inverse of the integer
sequence, and the multiplicative reciprocal sequence.
\end{abstract}

\begin{mainresultbox}
With
\[
\rho=0.33832185689920769519611262571701705318377460753296779557\ldots,
\]
\[
\alpha=\rho^{-1}=2.95576528565199497471481752412319458837549230466359659535\ldots,
\]
and
\[
 C=0.43992401257102530404090339143454476479808540794011985765\ldots,
\]
the first normalized terms are
\[
\begin{aligned}
 a_n\sim C\alpha^n n^{-3/2}\Biggl(&
 1+\frac{0.1004034800964014348}{n}
 +\frac{0.5039343151023035331}{n^2}\\[-1mm]
 &+\frac{1.972284908382278469}{n^3}
 +\frac{10.51734090413157399}{n^4}+\cdots\Biggr).
\end{aligned}
\]
All displayed decimals are consequences of the exact coefficient formulae,
not definitions of the coefficients.
\end{mainresultbox}

\tableofcontents
\clearpage

\section{Introduction and statement of scope}

\subsection{The sequence and its three meanings of ``inverse''}

The sequence considered here is
\[
 0,1,1,2,4,9,20,48,115,286,719,1842,4766,\ldots,
\]
with $a_0=0$ and $a_n$ equal to the number of isomorphism classes of rooted,
unlabelled, non-plane trees on $n$ vertices.  P\'olya's multiset construction
makes these objects especially natural in enumerative combinatorics.  The
same trees index elementary differentials and coefficients in Butcher
$B$-series, which explains the name \emph{Butcher tree sequence} used in
numerical analysis \cite{Butcher1972,McLachlan2017,OEISA000081}.

The word \emph{inverse} can refer to three different constructions, and a
correct asymptotic treatment must separate them.
\begin{enumerate}
\item The \emph{smooth functional inverse} $x=\nu(y)$ solves an asymptotic
      interpolation $\mathfrak a(x)=y$ for $x\to+\infty$.  This is the main
      inverse developed in this article.
\item The \emph{discrete generalized inverse}
      \[
       N(y):=\min\{n\ge1:a_n\ge y\}
      \]
      is a staircase.  Its bounded lattice oscillation cannot be represented
      by an ordinary descending power series; nevertheless its complete
      unbounded and vanishing asymptotic content is inherited from $\nu(y)$.
\item The \emph{multiplicative reciprocal} $1/a_n$ is not a functional
      inverse.  Its all-orders expansion is included in \cref{sec:reciprocal}
      because it follows at essentially no extra cost from the same
      coefficient calculus.
\end{enumerate}

\subsection{What ``full transseries'' means here}

There are two distinct levels of completeness.

First, the dominant singularity at $z=\rho$ gives a rigorous Poincar\'e
expansion to arbitrary algebraic order,
\[
 a_n=\frac{\rho^{-n}}{\sqrt\pi\,n^{3/2}}
 \left(\sum_{m=0}^{M-1}\frac{\tau_m}{n^m}
       +\Ord(n^{-M})\right)
 \qquad(M\ge1\text{ fixed}).
\]
The main algebraic contribution of this article is a simpler coefficient
formula for $\tau_m$ than a direct repeated singularity-analysis calculation:
it is a finite convolution of the odd Puiseux coefficients with universal
gamma-ratio coefficients.

Second, analytic continuation of the recursive generating equation produces
further singular germs.  Each germ gives a different exponential weight
$\sigma^{-n}$, so the coefficient sequence has a multi-sector expansion
ordered by $|\sigma|$.  The inherited family
$\sigma^m=\rho$ produces the characteristic scales
$\rho^{-n/m}$ with root-of-unity phases.  Additional solutions of the
critical equation $\zeta(z)=1/\e$ may also occur on continued sheets.  We
give a complete local coefficient recursion and a finite-radius transfer
theorem under explicit continuation and separation hypotheses.  This is the
appropriate rigorous meaning of a beyond-all-orders transseries: every
finite collection of accessible singularities can be expanded and
transferred, while a global assertion that the recursively generated list is
exhaustive would require a separate theorem on the full Riemann surface.

\begin{warningbox}
The dominant all-orders expansion and its algebraic functional inverse are
unconditional within standard singularity-analysis hypotheses known for
P\'olya trees.  The global exponential completion is stated as a finite-level
theorem under explicit analytic-continuation hypotheses.  This distinction
prevents a formal singularity grammar from being mistaken for an unproved
global classification of every sheet and every critical coincidence.
\end{warningbox}

\subsection{Relation to the supplied coefficient and transseries calculi}

The derivations are organized in the style of the supplied ProveIt drafts:
ordinary Bell polynomials encode powers of a series, complete exponential
Bell polynomials encode exponentials and Fa\`a di Bruno composition,
Bernoulli polynomials encode gamma-ratio transfer, and inversion is performed
only after the dominant carrier has been normalized
\cite{Comtet1974,ProveItCoefficientCalculus,ProveItTransseries}.  Formal identities are
kept separate from analytic remainder statements.  This makes the output
suitable both for symbolic implementation and for later formalization.

\subsection{Road map}

\Cref{sec:bell} fixes the coefficient conventions.  \Cref{sec:exact}
derives the exact generating equations and recurrence.  The Cayley normal
form and the universal square-root coefficients are developed in
\cref{sec:cayley}.  P\'olya-specific Puiseux coefficients are obtained in
\cref{sec:puiseux}.  \Cref{sec:transfer} proves the master formula for
$\tau_m$ and records numerical coefficients.  The logarithmic normal form,
reciprocal sequence, and ratio consequences are in
\cref{sec:lognormal,sec:reciprocal}.  The Lambert carrier and the two complete
forms of the functional inverse occupy \cref{sec:inverse-carrier,sec:inverse-pl}.
The discrete inverse is treated in \cref{sec:discrete}.  The singularity
transseries and its inverse are developed in
\cref{sec:global-transseries,sec:inverse-exponential}.  Algorithms,
validation, and reproducibility checks conclude the article.

\section{Coefficient calculi and notation}\label{sec:bell}

\subsection{Ordinary partial Bell polynomials}

\begin{definition}[Ordinary partial Bell polynomials]
For indeterminates $x_1,x_2,\ldots$, define
\[
 \left(\sum_{j\ge1}x_j u^j\right)^k
   =\sum_{n\ge k}\Bhat_{n,k}(x_1,x_2,\ldots)u^n,
\]
with $\Bhat_{0,0}=1$ and $\Bhat_{n,0}=0$ for $n>0$.  Equivalently,
\begin{equation}\label{eq:ordinary-bell-definition}
 \Bhat_{n,k}(x_1,x_2,\ldots)
 =\sum_{\substack{j_1+j_2+\cdots=k\\
                   j_1+2j_2+3j_3+\cdots=n}}
   \frac{k!}{j_1!j_2!\cdots}\prod_{m\ge1}x_m^{j_m}.
\end{equation}
\end{definition}

Thus, for arbitrary $\gamma$ and $U(u)=\sum_{j\ge1}x_ju^j$,
\begin{equation}\label{eq:generalized-binomial-bell}
 \Coeff{u}{n}(1+U(u))^\gamma
 =\sum_{k=0}^{n}\binom{\gamma}{k}
   \Bhat_{n,k}(x_1,x_2,\ldots).
\end{equation}
This single identity is the workhorse for fractional powers in the Puiseux
calculus and for negative powers in inverse-series calculations.

\subsection{Exponential Bell polynomials}

\begin{definition}[Partial and complete exponential Bell polynomials]
The partial exponential Bell polynomial $B_{n,k}$ is defined by
\[
 \frac1{k!}\left(\sum_{j\ge1}x_j\frac{u^j}{j!}\right)^k
 =\sum_{n\ge k}B_{n,k}(x_1,x_2,\ldots)\frac{u^n}{n!}.
\]
The complete exponential Bell polynomial is
\[
 Y_n(x_1,\ldots,x_n)=\sum_{k=0}^{n}B_{n,k}(x_1,\ldots,x_{n-k+1}).
\]
It is characterized by
\begin{equation}\label{eq:complete-bell-exp}
 \exp\!\left(\sum_{j\ge1}x_j\frac{u^j}{j!}\right)
 =\sum_{n\ge0}Y_n(x_1,\ldots,x_n)\frac{u^n}{n!}.
\end{equation}
\end{definition}

The Fa\`a di Bruno formula is then
\begin{equation}\label{eq:faa-di-bruno}
 \frac{\dd^n}{\dd z^n}f(g(z))
 =\sum_{k=1}^{n}f^{(k)}(g(z))
 B_{n,k}\bigl(g'(z),g''(z),\ldots,g^{(n-k+1)}(z)\bigr).
\end{equation}

\subsection{Bernoulli polynomials and gamma-ratio coefficients}

Let the Bernoulli polynomials be defined by
\[
 \frac{te^{xt}}{e^t-1}=\sum_{n\ge0}B_n(x)\frac{t^n}{n!}.
\]
For $\beta\in\C$ define
\begin{equation}\label{eq:kappa-def}
 \kappa_m(\beta):=
 \frac{(-1)^{m+1}}{m(m+1)}
 \bigl(B_{m+1}(-\beta)-B_{m+1}(1)\bigr),\qquad m\ge1,
\end{equation}
and define $g_r(\beta)$ by
\begin{equation}\label{eq:g-def-exp}
 \exp\!\left(\sum_{m\ge1}\kappa_m(\beta)u^m\right)
 =\sum_{r\ge0}g_r(\beta)u^r.
\end{equation}
Thus
\begin{equation}\label{eq:g-bell}
 g_r(\beta)=\frac1{r!}
 Y_r\bigl(1!\kappa_1(\beta),2!\kappa_2(\beta),\ldots,
           r!\kappa_r(\beta)\bigr),
\end{equation}
or, computationally,
\begin{equation}\label{eq:g-recurrence}
 g_0(\beta)=1,\qquad
 g_r(\beta)=\frac1r\sum_{m=1}^{r}m\kappa_m(\beta)g_{r-m}(\beta).
\end{equation}
The first values are
\begin{align*}
 g_0(\beta)&=1,\\
 g_1(\beta)&=\frac{\beta(\beta+1)}2,\\
 g_2(\beta)&=\frac{\beta(\beta+1)(\beta+2)(3\beta+1)}{24},\\
 g_3(\beta)&=\frac{\beta^2(\beta+1)^2(\beta+2)(\beta+3)}{48},\\
 g_4(\beta)&=\frac{\beta(\beta+1)(\beta+2)(\beta+3)(\beta+4)
 (15\beta^3+30\beta^2+5\beta-2)}{5760}.
\end{align*}

\begin{lemma}[Gamma-ratio expansion]\label{lem:gamma-ratio}
For fixed $\beta$ and $n\to\infty$ in a sector avoiding the negative real
axis,
\begin{equation}\label{eq:gamma-ratio}
 \frac{\Gamma(n-\beta)}{\Gamma(n+1)}
 \sim n^{-\beta-1}\sum_{r\ge0}\frac{g_r(\beta)}{n^r}.
\end{equation}
\end{lemma}

\begin{proof}
The standard Bernoulli-polynomial expansion of the logarithm of a gamma
ratio gives
\[
 \log\frac{\Gamma(n-\beta)}{\Gamma(n+1)}
 \sim-(\beta+1)\log n+
 \sum_{m\ge1}\frac{\kappa_m(\beta)}{n^m}.
\]
Exponentiating and applying \eqref{eq:complete-bell-exp} gives
\eqref{eq:gamma-ratio}, with the coefficient descriptions
\eqref{eq:g-bell} and \eqref{eq:g-recurrence}.
\end{proof}

\begin{summarybox}
There are only three coefficient mechanisms in the dominant calculation:
$\Bhat_{n,k}$ for powers of ordinary series, $B_{n,k}$ and $Y_n$ for
composition and exponentiation, and $g_r(\beta)$ for exact singularity
transfer.  Every later formula reduces to these finite polynomial objects.
\end{summarybox}

\section{Exact combinatorial and generating-function equations}\label{sec:exact}

\subsection{P\'olya's multiset equation}

A rooted non-plane tree consists of a root vertex carrying an unordered
multiset of rooted non-plane trees.  The unlabelled multiset construction and
P\'olya cycle-index substitution therefore give
\cite{Polya1937,FlajoletSedgewick2009,OEISA000081}
\begin{equation}\label{eq:polya-functional}
 T(z)=z\exp\!\left(\sum_{k\ge1}\frac{T(z^k)}{k}\right),
 \qquad T(z)=\sum_{n\ge1}a_nz^n.
\end{equation}
Separating the $k=1$ term yields
\begin{equation}\label{eq:zeta-normal}
 T(z)=\zeta(z)e^{T(z)},\qquad
 \zeta(z):=z\exp R(z),\qquad
 R(z):=\sum_{k\ge2}\frac{T(z^k)}{k}.
\end{equation}
The Euler-product form follows by expanding the logarithm:
\begin{equation}\label{eq:euler-product}
 T(z)=\frac{z}{\displaystyle\prod_{m\ge1}(1-z^m)^{a_m}}.
\end{equation}

\subsection{The exact quadratic-time recurrence}

Define the divisor transform
\begin{equation}\label{eq:divisor-transform}
 b_k:=\sum_{d\mid k}d\,a_d.
\end{equation}
Logarithmic differentiation of \eqref{eq:euler-product} gives
\[
 \frac{zT'(z)}{T(z)}=1+\sum_{k\ge1}b_kz^k.
\]
Coefficient extraction proves the following recurrence.

\begin{proposition}[Exact recurrence]\label{prop:exact-recurrence}
With $a_0=0$, $a_1=1$, one has, for $n\ge2$,
\begin{equation}\label{eq:exact-recurrence}
 a_n=\frac1{n-1}\sum_{k=1}^{n-1}
 \left(\sum_{d\mid k}d\,a_d\right)a_{n-k}
 =\frac1{n-1}\sum_{k=1}^{n-1}b_ka_{n-k}.
\end{equation}
In particular, the first $N$ values require $\Ord(N^2)$ arithmetic
operations when the divisor transform is updated by a sieve.
\end{proposition}

\begin{proof}
Multiplying the logarithmic-derivative identity by $T(z)$ gives
\[
 zT'(z)=T(z)+T(z)\sum_{k\ge1}b_kz^k.
\]
The coefficient of $z^n$ is
$n a_n=a_n+\sum_{k=1}^{n-1}b_ka_{n-k}$, which is
\eqref{eq:exact-recurrence}.  For the complexity statement, after computing
$a_d$ add $d a_d$ to $b_d,b_{2d},b_{3d},\ldots,b_{\lfloor N/d\rfloor d}$.
The sieve costs $\Ord(N\log N)$ additions, and the convolution in
\eqref{eq:exact-recurrence} costs $\Ord(N^2)$ arithmetic operations.
\end{proof}

\subsection{Critical characterization of the dominant singularity}

Let $\rho$ be the radius of convergence of $T$.  Since the coefficients are
nonnegative, Pringsheim's theorem puts a singularity at the positive point
$z=\rho$.  For $|z|<\sqrt\rho$, every $z^k$ with $k\ge2$ lies in the disc of
convergence of $T$, so $R$ and $\zeta$ are analytic there.  The implicit
equation $F(z,w)=w-\zeta(z)e^w=0$ becomes singular when
\[
 F=0,\qquad F_w=1-\zeta(z)e^w=0.
\]
Consequently the positive critical point satisfies
\begin{equation}\label{eq:critical-system}
 T(\rho)=1,\qquad \zeta(\rho)=e^{-1},
\end{equation}
or equivalently
\begin{equation}\label{eq:rho-scalar}
 \log\rho+1+\sum_{k\ge2}\frac{T(\rho^k)}{k}=0.
\end{equation}
The left side is strictly increasing on the relevant positive interval, so
this equation characterizes $\rho$ uniquely
\cite{Otter1948,FlajoletSedgewick2009,Genitrini2016}.

\begin{proposition}[Numerical constants]\label{prop:numerical-constants}
Solving \eqref{eq:rho-scalar} with the exact recurrence
\eqref{eq:exact-recurrence} gives
\begin{align}
 \rho&=0.33832185689920769519611262571701705318377460753296779557\ldots,
 \label{eq:rho-numeric}\\
 \alpha:=\rho^{-1}
 &=2.95576528565199497471481752412319458837549230466359659535\ldots,
 \label{eq:alpha-numeric}\\
 \lambda:=\log\alpha
 &=1.08375759724462466048271198089617427904966212420267223491\ldots.
 \label{eq:lambda-numeric}
\end{align}
\end{proposition}

\begin{proof}[Numerical certification strategy]
For any trial $r$ near $\rho$, every argument $r^k$, $k\ge2$, lies far
inside the disc $|z|<\rho$.  Hence
$T(r^k)=\sum_{n\ge1}a_nr^{kn}$ converges geometrically, and the omitted tail
can be bounded by a geometric majorant after finitely many exact
coefficients have been generated.  Interval evaluation of
\eqref{eq:rho-scalar} and its positive derivative then encloses the unique
root.  The displayed digits agree with the established Otter growth
constant $\alpha$ \cite{OEISOtterConstants}.
\end{proof}

\section{Cayley normal form and the universal square-root germ}\label{sec:cayley}

\subsection{The Cayley tree function}

Let $\Cay(w)$ be the solution analytic at the origin of
\begin{equation}\label{eq:cayley-def}
 \Cay(w)=w e^{\Cay(w)},\qquad \Cay(0)=0.
\end{equation}
Equivalently,
\begin{equation}\label{eq:cayley-lambert}
 \Cay(w)=-W_0(-w).
\end{equation}
The first critical point is $(w,\Cay)=(e^{-1},1)$.  By
\eqref{eq:zeta-normal}, the P\'olya-tree generating function is exactly
\begin{equation}\label{eq:T-C-zeta}
 T(z)=\Cay(\zeta(z))
\end{equation}
throughout the common domain of the branches selected at the origin.

To obtain the complete local germ of $\Cay$, put
\begin{equation}\label{eq:cayley-local-variables}
 w=\frac{1-u}{e},\qquad \Cay(w)=1-v.
\end{equation}
Substitution in \eqref{eq:cayley-def} gives
\begin{equation}\label{eq:u-v-relation}
 u=1-(1-v)e^v=\frac{v^2}{2}\Phi(v),
\end{equation}
where
\begin{equation}\label{eq:Phi-def}
 \Phi(v):=\frac{2\bigl(1-(1-v)e^v\bigr)}{v^2}
 =\sum_{m\ge0}\frac{2(m+1)}{(m+2)!}v^m
 =1+\frac23v+\frac14v^2+\cdots.
\end{equation}

\begin{theorem}[Universal Cayley Puiseux coefficients]\label{thm:cayley-coefficients}
On the branch reached from $0<w<e^{-1}$,
\begin{equation}\label{eq:cayley-puiseux}
 \Cay\!\left(\frac{1-u}{e}\right)
 =1+\sum_{r\ge1}c_r u^{r/2},
\end{equation}
where, for every $r\ge1$,
\begin{equation}\label{eq:c-r-lagrange}
 \boxed{
 c_r=-\frac{2^{r/2}}{r}\Coeff{v}{r-1}\Phi(v)^{-r/2}.}
\end{equation}
Equivalently, writing
$\phi_m=2(m+1)/(m+2)!$ for $m\ge1$,
\begin{equation}\label{eq:c-r-bell}
 \boxed{
 c_r=-\frac{2^{r/2}}r
 \sum_{k=0}^{r-1}\binom{-r/2}{k}
 \Bhat_{r-1,k}(\phi_1,\phi_2,\ldots).}
\end{equation}
\end{theorem}

\begin{proof}
Set $s=u^{1/2}$ with the positive square root for $u>0$.  Equation
\eqref{eq:u-v-relation} is equivalent to
\[
 v=s\,\Psi(v),\qquad \Psi(v)=\sqrt2\,\Phi(v)^{-1/2},
\]
with $\Psi(0)=\sqrt2\ne0$.  The Lagrange--B\"urmann formula gives
\[
 \Coeff{s}{r}v(s)=\frac1r\Coeff{v}{r-1}\Psi(v)^r
 =\frac{2^{r/2}}r\Coeff{v}{r-1}\Phi(v)^{-r/2}.
\]
Since $\Cay=1-v$, the coefficient in \eqref{eq:cayley-puiseux} is the
negative of this number, proving \eqref{eq:c-r-lagrange}.  Formula
\eqref{eq:c-r-bell} is \eqref{eq:generalized-binomial-bell} applied to
$\Phi(v)=1+\sum_{m\ge1}\phi_mv^m$.
\end{proof}

The first coefficients are
\begin{align}\label{eq:cayley-first-coeffs}
 c_1&=-\sqrt2,&
 c_2&=\frac23,&
 c_3&=-\frac{11\sqrt2}{36},&
 c_4&=\frac{43}{135},\notag\\
 c_5&=-\frac{769\sqrt2}{4320},&
 c_6&=\frac{1768}{8505},&
 c_7&=-\frac{680863\sqrt2}{5443200},&
 c_8&=\frac{3926}{25515}.
\end{align}
These constants are universal: the particular tree family enters only
through the analytic change of variable $u=1-e\zeta(z)$.

\subsection{Why the square root is unavoidable}

The critical equations for $F(w,y)=y-we^y$ are
$F(e^{-1},1)=F_y(e^{-1},1)=0$, whereas
$F_{yy}(e^{-1},1)=-1\ne0$ and $F_w(e^{-1},1)=-e\ne0$.
The Weierstrass preparation theorem therefore gives a quadratic local
branching.  The Lagrange formula above does more than establish the exponent:
it gives every coefficient in the branch germ by one finite coefficient
extraction.

\begin{checkpointbox}
All family dependence is now isolated in $q(s):=1-e\zeta(\rho(1-s))$.
The map $q\mapsto T$ is the universal composition
$1+\sum_{r\ge1}c_rq^{r/2}$.  This carrier-first normalization is the key
simplification behind the later all-orders coefficient formula.
\end{checkpointbox}

\section{All Puiseux coefficients at the Otter singularity}\label{sec:puiseux}

\subsection{The analytic perturbation series}

Put
\begin{equation}\label{eq:s-local}
 s:=1-\frac{z}{\rho}
\end{equation}
and define the analytic series
\begin{equation}\label{eq:q-def}
 q(s):=1-e\zeta\bigl(\rho(1-s)\bigr)
      =\sum_{m\ge1}q_ms^m.
\end{equation}
The constant term vanishes by \eqref{eq:critical-system}.  Taylor's formula
gives
\begin{equation}\label{eq:q-derivative}
 \boxed{
 q_m=\frac{(-1)^{m+1}e\rho^m}{m!}\zeta^{(m)}(\rho),\qquad m\ge1.}
\end{equation}
In particular $q_1=e\rho\zeta'(\rho)>0$.

Factor
\begin{equation}\label{eq:q-factor}
 q(s)=q_1s\bigl(1+U(s)\bigr),\qquad
 U(s)=\sum_{j\ge1}u_js^j,\qquad u_j:=\frac{q_{j+1}}{q_1}.
\end{equation}
The desired local expansion is
\begin{equation}\label{eq:T-puiseux-def}
 T(z)=1+\sum_{n\ge1}t_n\left(1-\frac z\rho\right)^{n/2}
     =1+\sum_{n\ge1}t_ns^{n/2}.
\end{equation}

\begin{theorem}[Finite Bell formula for every $t_n$]\label{thm:t-n-formula}
Let $c_r$ be given by \eqref{eq:c-r-lagrange}.  Then
\begin{equation}\label{eq:t-n-coeff}
 \boxed{
 t_n=\sum_{\substack{1\le r\le n\\r\equiv n\pmod 2}}
 c_r q_1^{r/2}
 \Coeff{s}{(n-r)/2}(1+U(s))^{r/2}.}
\end{equation}
Equivalently,
\begin{equation}\label{eq:t-n-bell}
 \boxed{
 t_n=\sum_{\substack{1\le r\le n\\r\equiv n\pmod 2}}
 c_rq_1^{r/2}
 \sum_{k=0}^{(n-r)/2}\binom{r/2}{k}
 \Bhat_{(n-r)/2,k}(u_1,u_2,\ldots).}
\end{equation}
Thus $t_n$ is a finite Laurent-polynomial expression in $q_1^{1/2}$,
polynomial in $q_2,\ldots,q_{(n+1)/2}$, with coefficients generated by
the universal rational--radical constants $c_1,\ldots,c_n$.
\end{theorem}

\begin{proof}
By \eqref{eq:T-C-zeta}, \eqref{eq:cayley-puiseux}, and
\eqref{eq:q-def},
\[
 T(z)=1+\sum_{r\ge1}c_rq(s)^{r/2}
 =1+\sum_{r\ge1}c_rq_1^{r/2}s^{r/2}(1+U(s))^{r/2}.
\]
A contribution to $s^{n/2}$ is possible exactly when $n-r$ is a
nonnegative even integer.  Extracting the remaining integral power of $s$
gives \eqref{eq:t-n-coeff}; inserting
\eqref{eq:generalized-binomial-bell} gives \eqref{eq:t-n-bell}.
\end{proof}

The first three nontrivial consequences are
\begin{align}
 t_1&=-\sqrt{2q_1},\label{eq:t1}\\
 t_2&=\frac23q_1,\label{eq:t2}\\
 t_3&=-\frac{11\sqrt2}{36}q_1^{3/2}
      -\frac{\sqrt2}{2}q_2q_1^{-1/2}.
\label{eq:t3}
\end{align}
Formula \eqref{eq:t-n-bell} is preferable to a rapidly growing list of
hand-expanded expressions because it preserves the triangular dependence and
is directly implementable.

\subsection{Derivatives of $\zeta$ from subcritical data}

The remaining task is to compute the derivatives in
\eqref{eq:q-derivative}.  Let
\begin{equation}\label{eq:L-def}
 L(z):=\log\zeta(z)=\log z+R(z),\qquad
 R(z)=\sum_{k\ge2}\frac{T(z^k)}k.
\end{equation}
Set
\begin{equation}\label{eq:ell-m-def}
 \ell_m:=L^{(m)}(\rho)
 =\frac{(-1)^{m-1}(m-1)!}{\rho^m}+R^{(m)}(\rho).
\end{equation}
Then the complete Bell polynomial gives
\begin{equation}\label{eq:zeta-derivative-bell}
 \boxed{
 \zeta^{(m)}(\rho)=e^{-1}Y_m(\ell_1,\ell_2,\ldots,\ell_m).}
\end{equation}

For $g_k(z)=z^k$, Fa\`a di Bruno yields
\begin{equation}\label{eq:R-derivative}
 \boxed{
 R^{(m)}(\rho)=
 \sum_{k\ge2}\frac1k\sum_{r=1}^{m}
 T^{(r)}(\rho^k)
 B_{m,r}\bigl(\fall{k}{1}\rho^{k-1},
               \fall{k}{2}\rho^{k-2},\ldots,
               \fall{k}{m-r+1}\rho^{k-m+r-1}\bigr).}
\end{equation}
Here $\fall{k}{j}=k(k-1)\cdots(k-j+1)$, with the convention that it is
zero for $j>k$.  Every argument $\rho^k$, $k\ge2$, is strictly subcritical,
so
\begin{equation}\label{eq:T-derivative-subcritical}
 T^{(r)}(\rho^k)=\sum_{n\ge r}\fall{n}{r}a_n\rho^{k(n-r)}
\end{equation}
converges geometrically.  Equations
\eqref{eq:q-derivative}, \eqref{eq:zeta-derivative-bell}, and
\eqref{eq:R-derivative} therefore constitute a complete and numerically
stable all-orders prescription for the input $q_m$.

\begin{remark}[A direct coefficient alternative]
For numerical work it is often faster to expand $R(\rho(1-s))$ directly.
Since
\[
 T\bigl(\rho^k(1-s)^k\bigr)
 =\sum_{n\ge1}a_n\rho^{kn}(1-s)^{kn},
\]
one has
\begin{equation}\label{eq:R-direct-coeff}
 \Coeff{s}{m}R\bigl(\rho(1-s)\bigr)
 =(-1)^m\sum_{k\ge2}\frac1k\sum_{n\ge1}
 a_n\rho^{kn}\binom{kn}{m}.
\end{equation}
The double sum converges geometrically.  Exponentiating this ordinary power
series by complete Bell polynomials recovers the same $q_m$ without explicit
derivatives.
\end{remark}

\subsection{Numerical Puiseux coefficients}

The first coefficients obtained from the preceding formulas are shown in
\cref{tab:t-coefficients}; they agree with the independent table in
\cite{Genitrini2016}.  Even-indexed coefficients belong to the analytic
part of the local germ; odd-indexed coefficients create the algebraic
coefficient asymptotics.

\begin{longtable}{@{}r@{\quad}r@{\qquad}r@{\quad}r@{}}
\caption{Puiseux coefficients in
$T(z)=1+\sum_{n\ge1}t_n(1-z/\rho)^{n/2}$.}
\label{tab:t-coefficients}\\
\toprule
$n$ & $t_n$ & $n$ & $t_n$\\
\midrule
\endfirsthead
\toprule
$n$ & $t_n$ & $n$ & $t_n$\\
\midrule
\endhead
1  & $-1.5594900203746408855$ & 10 & $\phantom{-}0.2059848312779074187$\\
2  & $\phantom{-}0.8106697078826992814$ & 11 & $-0.05272470849819056232$\\
3  & $-0.2854870216128456053$ & 12 & $\phantom{-}0.01702656875495366945$\\
4  & $\phantom{-}0.1653723657120838935$ & 13 & $-0.1523706243663253963$\\
5  & $-0.3424599704021542010$ & 14 & $\phantom{-}0.1737028832998504629$\\
6  & $\phantom{-}0.3174072259465285636$ & 15 & $-0.01447370373952704486$\\
7  & $-0.1077788002916310092$ & 16 & $-0.02189951761121556223$\\
8  & $\phantom{-}0.06138495705583510469$ & 17 & $-0.1445471935709097055$\\
9  & $-0.1952123835975564636$ & 18 & $\phantom{-}0.1760771088850177791$\\
\bottomrule
\end{longtable}

\begin{summarybox}
The universal coefficients $c_r$ are produced once by
\eqref{eq:c-r-lagrange}.  The family-specific derivatives of $\zeta$ are
produced by \eqref{eq:zeta-derivative-bell}--\eqref{eq:R-derivative}.  Their
composition is the finite Bell sum \eqref{eq:t-n-bell}.  No coefficient
requires ad hoc reversion or numerical fitting.
\end{summarybox}

\section{Coefficient transfer and the all-orders dominant expansion}\label{sec:transfer}

\subsection{Exact transfer of one singular monomial}

For $\beta\notin\Nzero$, the binomial identity gives the exact coefficient
formula
\begin{equation}\label{eq:exact-transfer}
 \Coeff{z}{n}\left(1-\frac z\rho\right)^\beta
 =\rho^{-n}\frac{\Gamma(n-\beta)}{\Gamma(-\beta)\Gamma(n+1)}.
\end{equation}
Combining \eqref{eq:exact-transfer} with \cref{lem:gamma-ratio} yields
\begin{equation}\label{eq:monomial-transfer-full}
 \Coeff{z}{n}\left(1-\frac z\rho\right)^\beta
 \sim\frac{\rho^{-n}n^{-\beta-1}}{\Gamma(-\beta)}
 \sum_{r\ge0}\frac{g_r(\beta)}{n^r}.
\end{equation}
For half-integers,
\begin{equation}\label{eq:half-gamma}
 \frac{1}{\Gamma(-j-\tfrac12)}
 =\frac{(-1)^{j+1}(2j+1)!!}{2^{j+1}\sqrt\pi},
 \qquad j\ge0.
\end{equation}

The even terms $t_{2j}(1-z/\rho)^j$ in the local expansion are analytic
polynomials and do not contribute to the large-$n$ coefficient tail.  The
odd terms contain all orders of the dominant algebraic asymptotics.

\subsection{The master coefficient theorem}

\begin{theorem}[Full dominant expansion]\label{thm:dominant-expansion}
For each fixed integer $M\ge1$,
\begin{equation}\label{eq:dominant-expansion-M}
 a_n=\frac{\rho^{-n}}{\sqrt\pi\,n^{3/2}}
 \left(\sum_{m=0}^{M-1}\frac{\tau_m}{n^m}
       +\Ord(n^{-M})\right),\qquad n\to\infty,
\end{equation}
where
\begin{equation}\label{eq:tau-master}
 \boxed{
 \tau_m=\sum_{j=0}^{m}
 D_j\,t_{2j+1}\,g_{m-j}\!\left(j+\tfrac12\right),
 \qquad
 D_j:=\frac{(-1)^{j+1}(2j+1)!!}{2^{j+1}}.}
\end{equation}
Equivalently,
\begin{equation}\label{eq:tau-master-gamma}
 \tau_m=\sqrt\pi\sum_{j=0}^{m}
 \frac{t_{2j+1}}{\Gamma(-j-\tfrac12)}
 g_{m-j}\!\left(j+\tfrac12\right).
\end{equation}
Together with \eqref{eq:c-r-bell}, \eqref{eq:t-n-bell},
\eqref{eq:q-derivative}, \eqref{eq:zeta-derivative-bell}, and
\eqref{eq:R-derivative}, this is a finite general formula for every
asymptotic coefficient $\tau_m$.
\end{theorem}

\begin{proof}
The P\'olya-tree generating function is analytic in a dented neighborhood
of its unique dominant singularity $\rho$, and its local expansion is
\eqref{eq:T-puiseux-def}.  Retain the odd terms
$t_{2j+1}(1-z/\rho)^{j+1/2}$ with $0\le j\le M-1$ and absorb the remaining
singular terms into an admissible remainder.  Applying
\eqref{eq:monomial-transfer-full} term by term gives
\[
 [z^n]t_{2j+1}\left(1-\frac z\rho\right)^{j+1/2}
 \sim t_{2j+1}\rho^{-n}n^{-j-3/2}
 \frac1{\Gamma(-j-\tfrac12)}
 \sum_{r\ge0}\frac{g_r(j+\tfrac12)}{n^r}.
\]
After factoring $\rho^{-n}/(\sqrt\pi n^{3/2})$, the coefficient of $n^{-m}$
is obtained from pairs $j+r=m$.  Inserting \eqref{eq:half-gamma} gives
\eqref{eq:tau-master}.  Standard transfer with a sufficiently long local
expansion gives the stated remainder for every fixed $M$.
\end{proof}

\begin{corollary}[Otter's leading constant]\label{cor:otter-constant}
The leading coefficient is
\begin{equation}\label{eq:tau0-C}
 \tau_0=-\frac{t_1}{2}=\sqrt{\frac{q_1}{2}},\qquad
 C:=\frac{\tau_0}{\sqrt\pi}
 =\sqrt{\frac{e\rho\zeta'(\rho)}{2\pi}}.
\end{equation}
Consequently
\begin{equation}\label{eq:leading-otter}
 a_n\sim C\alpha^n n^{-3/2},
\end{equation}
where $\alpha=\rho^{-1}$ and
\begin{equation}\label{eq:C-numeric}
 C=0.43992401257102530404090339143454476479808540794011985765\ldots.
\end{equation}
\end{corollary}

\begin{proof}
Set $m=0$ in \eqref{eq:tau-master}, use $D_0=-1/2$ and $g_0=1$, and then
insert $t_1=-\sqrt{2q_1}$ from \eqref{eq:t1} and
$q_1=e\rho\zeta'(\rho)$ from \eqref{eq:q-derivative}.
\end{proof}

\subsection{First coefficients in expanded form}

The first five instances of \eqref{eq:tau-master} are
\begin{align}
 \tau_0={}&-\frac{t_1}{2},\label{eq:tau0-explicit}\\
 \tau_1={}&-\frac{3(t_1-4t_3)}{16},\label{eq:tau1-explicit}\\
 \tau_2={}&-\frac{5(5t_1-72t_3+96t_5)}{256},\label{eq:tau2-explicit}\\
 \tau_3={}&-\frac{105(t_1-44t_3+160t_5-128t_7)}{2048},
 \label{eq:tau3-explicit}\\
 \tau_4={}&-\frac{21}{65536}
 \bigl(79t_1-10800t_3+81600t_5-161280t_7+92160t_9\bigr).
 \label{eq:tau4-explicit}
\end{align}
The compact convolution \eqref{eq:tau-master} scales much better than these
fully expanded expressions.

\subsection{Numerical coefficients}

\begin{longtable}{@{}r@{\quad}r@{\qquad}r@{\quad}r@{}}
\caption{Coefficients in
$a_n\sim\rho^{-n}(\pi n^3)^{-1/2}\sum_{m\ge0}\tau_mn^{-m}$.}
\label{tab:tau-coefficients}\\
\toprule
$m$ & $\tau_m$ & $m$ & $\tau_m$\\
\midrule
\endfirsthead
\toprule
$m$ & $\tau_m$ & $m$ & $\tau_m$\\
\midrule
\endhead
0  & $0.7797450101873204428$ & 10 & $1.599693249293173313\times10^{7}$\\
1  & $0.07828911261061096206$ & 11 & $2.928225313353392231\times10^{8}$\\
2  & $0.3929402676631860185$ & 12 & $5.914523441293932519\times10^{9}$\\
3  & $1.537879315978838091$ & 13 & $1.305991927898971332\times10^{11}$\\
4  & $8.200844090435596159$ & 14 & $3.128498399789518761\times10^{12}$\\
5  & $57.29291473494343788$ & 15 & $8.078305401468884289\times10^{13}$\\
6  & $503.0445050262735818$ & 16 & $2.236301680891633732\times10^{15}$\\
7  & $5359.600933884326033$ & 17 & $6.605960869699189564\times10^{16}$\\
8  & $67342.06920114653048$ & 18 & $2.073828085209893749\times10^{18}$\\
9  & $975425.4970695924732$ & & \\
\bottomrule
\end{longtable}

It is often preferable to factor out the leading constant.  Define
\begin{equation}\label{eq:S-def}
 S(u):=1+\sum_{m\ge1}s_mu^m,
 \qquad s_m:=\frac{\tau_m}{\tau_0}.
\end{equation}
Then
\begin{equation}\label{eq:normalized-forward}
 \boxed{
 a_n\sim C\alpha^n n^{-3/2}S(n^{-1}).}
\end{equation}
The first normalized coefficients are
\begin{align*}
 s_1&=0.1004034800964014347637,\\
 s_2&=0.5039343151023035331022,\\
 s_3&=1.972284908382278469206,\\
 s_4&=10.51734090413157399274,\\
 s_5&=73.47647498402047106417,\\
 s_6&=645.139755245661293\ldots,\\
 s_7&=6873.530274463408\ldots,\\
 s_8&=86364.21948371140\ldots.
\end{align*}

\subsection{Remainders, divergence, and optimal truncation}

The statement ``full expansion'' means that every fixed truncation has the
next Poincar\'e order; it does not assert convergence of
$\sum\tau_mn^{-m}$.  The rapid growth in \cref{tab:tau-coefficients} is
consistent with a divergent, Gevrey-type expansion controlled by more distant
singular structure.  At a fixed moderately sized $n$, adding terms first
improves the approximation and eventually worsens it.  This is the usual
least-term principle.

A useful diagnostic follows directly from the terms
$|\tau_m|n^{-m}$.  Truncate near the smallest available term, and estimate
the algebraic truncation uncertainty by the first omitted term.  A rigorous
bound requires uniform control of the local remainder; the least-term rule is
a numerical strategy rather than a substitute for that bound.

\section{Logarithmic normal form and related sequence transforms}\label{sec:lognormal}

\subsection{The logarithm of the asymptotic series}

Define
\begin{equation}\label{eq:H-def}
 H(u):=\log S(u)=\sum_{m\ge1}h_mu^m.
\end{equation}
By the ordinary Bell-polynomial convention,
\begin{equation}\label{eq:h-bell}
 \boxed{
 h_m=\sum_{k=1}^{m}\frac{(-1)^{k-1}}{k}
 \Bhat_{m,k}(s_1,s_2,\ldots).}
\end{equation}
Equivalently, differentiating $S=e^H$ gives the triangular recurrence
\begin{equation}\label{eq:h-recurrence}
 h_m=s_m-\frac1m\sum_{k=1}^{m-1}k h_k s_{m-k}.
\end{equation}
Therefore
\begin{equation}\label{eq:log-a-n}
 \boxed{
 \log a_n\sim n\lambda-\frac32\log n+\log C
 +\sum_{m\ge1}\frac{h_m}{n^m}.}
\end{equation}
The first values are
\begin{align*}
 h_1&=0.1004034800964014347637,\\
 h_2&=0.4988938856945692935703,\\
 h_3&=1.922025533841821821591,\\
 h_4&=10.19739642337376887053,\\
 h_5&=71.47146706207073604412,\\
 h_6&=630.859937607337\ldots.
\end{align*}
This logarithmic normal form is the natural input for functional inversion.

\subsection{Ratio expansion and eventual monotonicity}

Subtracting \eqref{eq:log-a-n} at $n$ and $n+1$ gives
\begin{align}
 \log\frac{a_{n+1}}{a_n}
 \sim{}&\lambda-\frac{3}{2n}
 +\frac{\frac34-h_1}{n^2}
 +\frac{-\frac12+h_1-2h_2}{n^3}+\cdots.
 \label{eq:ratio-log}
\end{align}
Exponentiation gives a complete expansion for $a_{n+1}/a_n$.  In particular,
\begin{equation}\label{eq:ratio-limit}
 \frac{a_{n+1}}{a_n}\longrightarrow\alpha>1,
\end{equation}
so the sequence is strictly increasing for all sufficiently large $n$, and
the generalized inverse in \cref{sec:discrete} is well defined without any
large-$n$ ambiguity.

\section{The multiplicative reciprocal sequence}\label{sec:reciprocal}

Although the main inverse is functional, the reciprocal sequence has an
immediate all-orders expansion.  Define
\begin{equation}\label{eq:reciprocal-series}
 R_{\!*}(u):=\frac1{S(u)}=\sum_{m\ge0}r_mu^m,
 \qquad r_0=1.
\end{equation}
Then
\begin{equation}\label{eq:reciprocal-asymptotic}
 \boxed{
 \frac1{a_n}\sim C^{-1}\rho^n n^{3/2}
 \sum_{m\ge0}\frac{r_m}{n^m}.}
\end{equation}
The coefficients have both a triangular and a Bell-polynomial form:
\begin{align}
 r_m&=-\sum_{k=1}^{m}s_kr_{m-k},\label{eq:r-recurrence}\\
 r_m&=\sum_{k=1}^{m}(-1)^k
 \Bhat_{m,k}(s_1,s_2,\ldots),\qquad m\ge1.
 \label{eq:r-bell}
\end{align}
For example,
\[
 r_1=-s_1,
 \qquad r_2=s_1^2-s_2,
 \qquad r_3=-s_1^3+2s_1s_2-s_3.
\]
This section also illustrates why it is essential not to call $1/a_n$ the
functional inverse: its exponential scale is $\rho^n$, whereas the
functional inverse grows only logarithmically in its argument.

\section{Functional inversion: the exact Lambert carrier}\label{sec:inverse-carrier}

\subsection{The smooth asymptotic model}

Set
\begin{equation}\label{eq:p-def}
 p:=\frac32
\end{equation}
and regard
\begin{equation}\label{eq:A-model}
 \mathcal A(x):=C e^{\lambda x}x^{-p}S(x^{-1})
\end{equation}
as a formal asymptotic interpolation of $a_n$.  For sufficiently large real
$x$, every finite truncation is strictly increasing.  We denote by
$\nu(y)$ the large real solution of
\begin{equation}\label{eq:inverse-def}
 \mathcal A(\nu(y))=y,
\end{equation}
interpreted asymptotically: substitution of an $M$-term approximation for
$S$ determines $\nu$ through the corresponding order.

Taking logarithms gives the canonical reversion problem
\begin{equation}\label{eq:inverse-log-equation}
 \log\frac yC=\lambda x-p\log x+H(x^{-1}).
\end{equation}
The slowly varying factor $H$ should not be expanded until the elementary
exponential--power carrier has been inverted exactly.

\subsection{Exact inversion of the exponential--power carrier}

Let $x_0(y)$ solve
\begin{equation}\label{eq:carrier-equation}
 y=C e^{\lambda x_0}x_0^{-p}.
\end{equation}
Raising both sides to the power $-1/p$ gives
\[
 -\frac{\lambda}{p}\left(\frac Cy\right)^{1/p}
 =-\frac{\lambda x_0}{p}e^{-\lambda x_0/p}.
\]
Hence
\begin{equation}\label{eq:carrier-W-general}
 \boxed{
 x_0(y)=-\frac p\lambda
 W_{-1}\!\left(-\frac\lambda p\left(\frac Cy\right)^{1/p}\right).}
\end{equation}
For the rooted-tree exponent $p=3/2$,
\begin{equation}\label{eq:carrier-W-tree}
 \boxed{
 x_0(y)=-\frac{3}{2\lambda}
 W_{-1}\!\left(-\frac{2\lambda}{3}
 \left(\frac Cy\right)^{2/3}\right).}
\end{equation}
The Lambert representation uses the standard real branches described in
\cite{Corless1996}.  The branch $W_{-1}$ is mandatory: its value tends to $-\infty$ as the
argument approaches $0$ from below, and therefore $x_0(y)\to+\infty$.
The principal branch $W_0$ would give the irrelevant small solution.

\subsection{All algebraic corrections around the carrier}

Write
\begin{equation}\label{eq:D-ansatz}
 \nu(y)=x_0(y)+D\bigl(x_0(y)^{-1}\bigr),\qquad
 D(t)=\sum_{m\ge1}d_mt^m.
\end{equation}
Subtracting the carrier equation from \eqref{eq:inverse-log-equation} gives
an exact formal correction equation.

\begin{theorem}[Lambert-carrier correction equation]\label{thm:D-equation}
The series $D(t)\in t\R[[t]]$ is the unique formal solution of
\begin{equation}\label{eq:D-equation}
 \boxed{
 \lambda D(t)-p\log\bigl(1+tD(t)\bigr)
 +H\!\left(\frac{t}{1+tD(t)}\right)=0.}
\end{equation}
Its coefficients are triangularly determined by
\begin{equation}\label{eq:d-coeff-extraction}
 \boxed{
 d_m=\Coeff{t}{m}\left\{
 \frac p\lambda\log\bigl(1+tD_{<m}(t)\bigr)
 -\frac1\lambda H\!\left(
 \frac{t}{1+tD_{<m}(t)}\right)\right\},}
\end{equation}
where $D_{<m}(t)=\sum_{j=1}^{m-1}d_jt^j$.
\end{theorem}

\begin{proof}
Put $x=x_0+D$ in \eqref{eq:inverse-log-equation} and subtract
$\lambda x_0-p\log x_0=\log(y/C)$.  Since $t=x_0^{-1}$,
\[
 0=\lambda D-p\log\left(1+\frac D{x_0}\right)
 +H\left(\frac1{x_0+D}\right),
\]
which is \eqref{eq:D-equation}.  Its linear coefficient in $D$ at $t=0$ is
$\lambda\ne0$, so the formal implicit-function theorem gives a unique
solution in $t\R[[t]]$.  Moreover, the coefficient of $t^m$ on the right
side of the rearranged equation
\[
 D=\frac p\lambda\log(1+tD)
   -\frac1\lambda H\left(\frac{t}{1+tD}\right)
\]
depends only on $d_1,\ldots,d_{m-1}$, proving the recursion.
\end{proof}

\subsection{Bell-polynomial form of the inverse coefficients}

Expanding the logarithm and the negative powers in
\eqref{eq:d-coeff-extraction} yields a fully finite formula.  With the
convention $\Bhat_{0,0}=1$ and impossible indices equal to zero,
\begin{align}
 d_m={}&\frac p\lambda
 \sum_{r=1}^{m}\frac{(-1)^{r-1}}r
 \Bhat_{m-r,r}(d_1,d_2,\ldots)
 \label{eq:d-bell-first}\\
 &-\frac1\lambda\sum_{k=1}^{m}h_k
 \sum_{r=0}^{m-k}(-1)^r\binom{k+r-1}{r}
 \Bhat_{m-k-r,r}(d_1,d_2,\ldots).
 \label{eq:d-bell-second}
\end{align}
Although $d_m$ appears in the notation on the right, the index constraints
ensure that only $d_1,\ldots,d_{m-1}$ occur.

The first coefficients are
\begin{align}
 d_1={}&-\frac{h_1}{\lambda},\label{eq:d1}\\
 d_2={}&-\frac{ph_1+\lambda h_2}{\lambda^2},\label{eq:d2}\\
 d_3={}&-\frac{\lambda h_1^2+p^2h_1+p\lambda h_2+\lambda^2h_3}
 {\lambda^3},\label{eq:d3}\\
 d_4={}&-\frac{1}{2\lambda^4}
 \bigl(5p\lambda h_1^2+6\lambda^2h_1h_2+2p^3h_1
 +2p^2\lambda h_2+2p\lambda^2h_3+2\lambda^3h_4\bigr).
 \label{eq:d4}
\end{align}
Numerically,
\begin{align*}
 d_1&=-0.09264385352561312984200,\\
 d_2&=-0.5885630399313472628118,\\
 d_3&=-2.596680167407531021207,\\
 d_4&=-13.14905339996706568529,\\
 d_5&=-85.49870151462556652201,\\
 d_6&=-711.262914832946\ldots,\\
 d_7&=-7306.78827028497\ldots,\\
 d_8&=-89561.7608826372\ldots.
\end{align*}
Therefore the Lambert-normalized inverse is
\begin{equation}\label{eq:inverse-lambert-series}
 \boxed{
 \nu(y)\sim x_0(y)+\frac{d_1}{x_0(y)}+
 \frac{d_2}{x_0(y)^2}+\frac{d_3}{x_0(y)^3}+\cdots.}
\end{equation}

\begin{theorem}[Asymptotic validity of the correction series]
\label{thm:inverse-validity}
For every fixed $M\ge1$, let
\[
 \nu_M(y):=x_0(y)+\sum_{m=1}^{M-1}d_mx_0(y)^{-m}.
\]
Then substitution into the logarithmic forward model gives
\begin{equation}\label{eq:inverse-residual}
 \lambda\nu_M-p\log\nu_M+H(\nu_M^{-1})
 -\log(y/C)=\Ord(x_0^{-M}).
\end{equation}
Since the derivative of the left side with respect to $x$ is
$\lambda+\Ord(x^{-1})$, the corresponding inverse error is
\begin{equation}\label{eq:inverse-error}
 \nu(y)-\nu_M(y)=\Ord(x_0(y)^{-M}).
\end{equation}
\end{theorem}

\begin{proof}
Equation \eqref{eq:D-equation} determines the coefficient of each power of
$t$ uniquely.  Truncation after $t^{M-1}$ leaves a residual
$\Ord(t^M)$.  The derivative of the logarithmic forward map is bounded away
from zero for large $x$, so the mean-value theorem converts the residual
into an inverse error of the same order.
\end{proof}

\section{The inverse in polynomial--logarithmic form}\label{sec:inverse-pl}

The Lambert form is compact and numerically stable.  For algebraic
manipulation it is useful to eliminate $W_{-1}$ and expand directly in the
canonical polynomial--logarithmic scale.

\subsection{Normalization}

Define
\begin{equation}\label{eq:X-ell-q}
 X:=\frac1\lambda\log\frac yC,
 \qquad \ell:=\log X,
 \qquad q:=\frac p\lambda=\frac{3}{2\lambda},
 \qquad \eta_k:=\frac{h_k}{\lambda}.
\end{equation}
Then \eqref{eq:inverse-log-equation} becomes
\begin{equation}\label{eq:normalized-inverse-equation}
 x-q\log x+G(x^{-1})=X,
 \qquad G(u):=\sum_{k\ge1}\eta_ku^k.
\end{equation}
Seek
\begin{equation}\label{eq:PL-ansatz}
 \boxed{
 \nu(y)\sim X+q\ell+\sum_{m\ge1}\frac{P_m(\ell)}{X^m}.}
\end{equation}

\begin{theorem}[Polynomial--log inverse recursion]
\label{thm:P-recursion}
Put $t=X^{-1}$ and
$R(t,\ell)=\sum_{m\ge1}P_m(\ell)t^m$.  Then $R$ is the unique formal
solution of
\begin{equation}\label{eq:R-PL-equation}
 \boxed{
 R-q\log\bigl(1+t(q\ell+R)\bigr)
 +G\!\left(\frac{t}{1+t(q\ell+R)}\right)=0.}
\end{equation}
Consequently
\begin{equation}\label{eq:P-coeff-recurrence}
 \boxed{
 P_m(\ell)=\Coeff{t}{m}\left\{
 q\log\bigl(1+t(q\ell+R_{<m})\bigr)
 -G\!\left(\frac{t}{1+t(q\ell+R_{<m})}\right)
 \right\},}
\end{equation}
where $R_{<m}=\sum_{j=1}^{m-1}P_j(\ell)t^j$.  Each $P_m$ is a polynomial
of degree $m$, with leading term
\begin{equation}\label{eq:P-leading}
 P_m(\ell)=\frac{(-1)^{m+1}}m q^{m+1}\ell^m+
 \text{a polynomial of degree at most $m-1$}.
\end{equation}
\end{theorem}

\begin{proof}
Insert $x=X+q\ell+R=X(1+t(q\ell+R))$ into
\eqref{eq:normalized-inverse-equation}.  Since
$\log x=\ell+\log(1+t(q\ell+R))$, the terms $q\ell$ cancel and one obtains
\eqref{eq:R-PL-equation}.  The same $t$-adic triangularity used in
\cref{thm:D-equation} proves existence, uniqueness, and
\eqref{eq:P-coeff-recurrence}.  The highest power of $\ell$ at order $t^m$
comes from the $m$th term in the logarithm of $1+tq\ell$ and equals
$(-1)^{m+1}q^{m+1}\ell^m/m$.
\end{proof}

\subsection{A completely explicit Bell recursion}

For $N,r\ge0$ define
\begin{equation}\label{eq:A-N-r}
 \mathcal P_{N,r}(\ell):=
 \Coeff{t}{N}\left(q\ell+\sum_{j\ge1}P_j(\ell)t^j\right)^r
 =\sum_{j=0}^{\min(r,N)}\binom rj(q\ell)^{r-j}
 \Bhat_{N,j}(P_1,P_2,\ldots).
\end{equation}
Then \eqref{eq:P-coeff-recurrence} is equivalent to
\begin{align}
 P_m(\ell)={}&q\sum_{r=1}^{m}\frac{(-1)^{r-1}}r
 \mathcal P_{m-r,r}(\ell)\label{eq:P-bell-1}\\
 &-\sum_{k=1}^{m}\eta_k\sum_{r=0}^{m-k}
 (-1)^r\binom{k+r-1}{r}\mathcal P_{m-k-r,r}(\ell).
 \label{eq:P-bell-2}
\end{align}
The index bounds again ensure that only $P_1,\ldots,P_{m-1}$ appear on the
right.  Thus \eqref{eq:P-bell-1}--\eqref{eq:P-bell-2} is a general finite
coefficient formula in exactly the ordinary-Bell-polynomial language used
for the forward expansion.

\subsection{First polynomial coefficients}

The first four polynomials are
\begin{align}
 P_1(\ell)={}&q^2\ell-\eta_1,\label{eq:P1}\\
 P_2(\ell)={}&-\frac12q^3\ell^2
 +(\eta_1q+q^3)\ell-\eta_1q-\eta_2,
 \label{eq:P2}\\
 P_3(\ell)={}&\frac13q^4\ell^3
 -\left(\eta_1q^2+\frac32q^4\right)\ell^2
 +\left(3\eta_1q^2+2\eta_2q+q^4\right)\ell
 \notag\\
 &-\eta_1^2-\eta_1q^2-\eta_2q-\eta_3,
 \label{eq:P3}\\
 P_4(\ell)={}&-\frac14q^5\ell^4
 +\left(\eta_1q^3+\frac{11}{6}q^5\right)\ell^3
 \notag\\
 &+\left(-\frac{11}{2}\eta_1q^3-3\eta_2q^2-3q^5\right)\ell^2
 \notag\\
 &+\left(3\eta_1^2q+6\eta_1q^3+5\eta_2q^2+3\eta_3q+q^5\right)\ell
 \notag\\
 &-\frac52\eta_1^2q-3\eta_1\eta_2-\eta_1q^3
 -\eta_2q^2-\eta_3q-\eta_4.
 \label{eq:P4}
\end{align}
For the rooted-tree constants,
\begin{align*}
 q&=1.38407334242790228719882463131\ldots,\\
 \eta_1&=0.09264385352561312984200\ldots,\\
 \eta_2&=0.46033715192675089877\ldots,\\
 \eta_3&=1.77348286990414903\ldots,\\
 \eta_4&=9.4092963678408456\ldots.
\end{align*}
In particular, the first two scales are
\begin{equation}\label{eq:inverse-leading-simple}
 \boxed{
 \nu(y)=\frac1\lambda\log\frac yC
 +\frac{3}{2\lambda}\log\!\left(
 \frac1\lambda\log\frac yC\right)
 +\Ord\!\left(\frac{\log\log y}{\log y}\right).}
\end{equation}
The complete correction at every order is supplied by
\cref{thm:P-recursion}.

\begin{checkpointbox}
The two inverse forms are equivalent but serve different purposes.
Equation \eqref{eq:inverse-lambert-series} keeps the entire logarithmic
carrier inside $W_{-1}$ and is numerically superior.  Equation
\eqref{eq:PL-ansatz} expands the same answer in the canonical field
$\R[\log X][[X^{-1}]]$ and exposes every coefficient as a polynomial.
\end{checkpointbox}

\section{The discrete generalized inverse and lattice effects}\label{sec:discrete}

For large $n$, $a_n$ is strictly increasing.  Define
\begin{equation}\label{eq:N-def}
 N(y):=\min\{n\ge1:a_n\ge y\}.
\end{equation}
The step width on the logarithmic $y$-axis tends to
$\log\alpha$, because $a_{n+1}/a_n\to\alpha$.  Hence no continuous
asymptotic series can determine the position inside each unit step.

\begin{theorem}[Smooth versus discrete inverse]\label{thm:discrete-inverse}
Let $\nu$ be any smooth asymptotic inverse whose forward interpolation
matches \eqref{eq:normalized-forward} to all algebraic orders.  Then
\begin{equation}\label{eq:N-nu-O1}
 N(y)=\nu(y)+\Ord(1),\qquad y\to\infty.
\end{equation}
Moreover, at sequence values,
\begin{equation}\label{eq:nu-an-superalg}
 \nu(a_n)=n+\Ord(n^{-M})
\end{equation}
for every fixed $M$, provided the forward interpolation is retained to the
corresponding order.
\end{theorem}

\begin{proof}
The logarithmic derivative of the dominant interpolation tends to
$\lambda>0$.  Multiplying $y$ by a bounded factor therefore shifts its
smooth inverse by a bounded amount, while consecutive exact values differ
by the asymptotically fixed factor $\alpha$.  This proves
\eqref{eq:N-nu-O1}.  At $y=a_n$, an $M$-order relative forward error is
$\Ord(n^{-M})$; division by the nonzero limiting logarithmic derivative
preserves that order in the inverse variable.
\end{proof}

A lattice-completed approximation may be written formally as
\begin{equation}\label{eq:lattice-completion}
 N_{\mathrm{lat}}(y):=\left\lceil\nu(y)\right\rceil.
\end{equation}
It is correct away from vanishing neighborhoods of the exact jump values,
but no asymptotic formula can choose the value exactly at a discontinuity
without specifying the convention in \eqref{eq:N-def}.  The bounded
sawtooth $\lceil x\rceil-x$ is therefore a genuine extra sector, not a term
of the descending power-logarithmic series.

\begin{warningbox}
Exponentially small changes in a smooth interpolation can move a value
across an integer boundary.  Thus the algebraic inverse transseries is
canonical, but an exact below-unit description of the staircase requires
both a chosen interpolation and an explicit rounding convention.
\end{warningbox}

\section{Beyond-all-orders coefficient transseries}\label{sec:global-transseries}

\subsection{Why one dominant power series is not the whole story}

The expansion in \cref{thm:dominant-expansion} is complete in the Poincar\'e
sense at the dominant exponential scale $\rho^{-n}$.  It cannot detect a
term such as $\sigma^{-n}$ with $|\sigma|>\rho$, because
\[
 \frac{|\sigma|^{-n}}{\rho^{-n}}=\left(\frac\rho{|\sigma|}\right)^n
\]
decays faster than every power of $n^{-1}$.  Such terms are precisely the
next transseries sectors.

The recursive equation itself predicts inherited singularities.  If $T$ has
a singularity at $\tau$, then the summand $T(z^k)/k$ in $R(z)$ has singular
germs at every $k$th root of $\tau$.  In addition, even when $R$ is regular,
the outer Cayley function becomes critical whenever
\begin{equation}\label{eq:continued-critical}
 \zeta(z)=e^{-1}
\end{equation}
on the chosen continued sheet.  Thus there are two mechanisms:
\begin{enumerate}
\item \emph{inherited preimages}: $z^k$ reaches an earlier singularity;
\item \emph{new Cayley critical points}: the continued analytic input reaches
      the branch value $e^{-1}$.
\end{enumerate}

The visible inherited skeleton generated by the principal singularity is
\begin{equation}\label{eq:root-skeleton}
 \mathscr R:=\left\{
 \rho^{1/m}e^{2\pi\ii j/m}:m\ge1,\ 0\le j<m
 \right\}.
\end{equation}
Its level-$m$ coefficient weights are
\begin{equation}\label{eq:root-sector-weight}
 \sigma_{m,j}^{-n}=\rho^{-n/m}e^{-2\pi\ii jn/m}.
\end{equation}
For integer $n$, the phase is periodic of period dividing $m$.  Conjugate
singularities combine into real oscillatory sectors.

\begin{warningbox}
Equation \eqref{eq:root-skeleton} is a guaranteed recursive skeleton of
candidate inherited germs, not by itself an exhaustive global singularity
theorem.  Analytic continuation may encounter additional critical solutions
of \eqref{eq:continued-critical}, and a critical point may coalesce with an
inherited germ.  The local calculus below handles all of these cases once
the relevant continued germs have been identified.
\end{warningbox}

\subsection{General transfer from a singular germ}

Let $\sigma\ne0$ be an accessible singularity on a specified sheet, and put
$u=1-z/\sigma$.  Suppose its nonanalytic local part has a generalized
Puiseux expansion
\begin{equation}\label{eq:general-local-germ}
 T(z)\sim A_\sigma(u)+
 \sum_{\beta\in E_\sigma}c_{\sigma,\beta}u^\beta,
\end{equation}
where $A_\sigma$ is analytic, $E_\sigma$ is locally finite and well ordered
by real part, and no displayed $\beta$ is a nonnegative integer.

\begin{theorem}[Universal singular-sector coefficient formula]
\label{thm:general-sector}
Under the usual dented-domain condition at $\sigma$, the contribution of
\eqref{eq:general-local-germ} to the coefficient transseries is
\begin{equation}\label{eq:general-sector}
 \boxed{
 \mathfrak S_\sigma(n)
 \sim \sigma^{-n}
 \sum_{\beta\in E_\sigma}
 \frac{c_{\sigma,\beta}}{\Gamma(-\beta)}
 n^{-\beta-1}
 \sum_{r\ge0}\frac{g_r(\beta)}{n^r}.}
\end{equation}
Thus the coefficient multiplying
$\sigma^{-n}n^{-\beta-1-r}$ is exactly
\begin{equation}\label{eq:general-sector-coeff}
 \boxed{
 A_{\sigma,\beta,r}
 =\frac{c_{\sigma,\beta}}{\Gamma(-\beta)}g_r(\beta).}
\end{equation}
\end{theorem}

\begin{proof}
Apply the exact transfer identity \eqref{eq:exact-transfer}, with $\rho$
replaced by $\sigma$, to every noninteger monomial in the local germ, and
then use the gamma-ratio expansion \eqref{eq:gamma-ratio}.  Analytic powers
belong to the regular local part and do not contribute an infinite
large-$n$ tail.  Hankel-contour deformation justifies summing any prescribed
finite truncation under the dented-domain hypothesis.
\end{proof}

This theorem is the full analogue of \eqref{eq:tau-master}; it allows
arbitrary ramification exponents, not only half-integers.

\subsection{A finite-radius transseries theorem}

Fix $R$ with $\rho<R<1$.  The following hypothesis is explicit enough for a
rigorous finite-level statement.

\begin{definition}[Continuation hypothesis $\mathsf H(R)$]
We say that $\mathsf H(R)$ holds if a chosen continuation of $T$ is analytic
in $|z|<R$ after removing finitely many nonintersecting cuts beginning at a
finite set $\Sigma_R$, every $\sigma\in\Sigma_R$ has a generalized Puiseux
germ of the form \eqref{eq:general-local-germ}, the germs are
Delta-regular, and the remainder of the deformed coefficient contour lies
on a circle of radius strictly larger than every retained singularity and
at least $R$ up to harmless indentations.
\end{definition}

\begin{theorem}[Finite-level coefficient transseries]\label{thm:finite-level}
Assume $\mathsf H(R)$.  Order $\Sigma_R$ by increasing modulus.  For any
finite truncation of the exponent and inverse-power sums,
\begin{equation}\label{eq:finite-level-transseries}
 a_n\sim\sum_{\sigma\in\Sigma_R}\mathfrak S_\sigma(n)
       +\Ord(R^{-n}n^K)
\end{equation}
for a finite exponent $K$ determined by the local remainder bounds.  Every
coefficient in every retained sector is given by
\eqref{eq:general-sector-coeff}.
\end{theorem}

\begin{proof}
Deform the Cauchy coefficient contour into the sum of small Hankel contours
around the cuts issuing from points in $\Sigma_R$ plus an outer remainder
contour.  Apply \cref{thm:general-sector} on each Hankel contour.  The outer
contour is exponentially bounded by $R^{-n}$ times a polynomial factor
coming from the allowed boundary growth.  Since only finitely many
singularities and local terms are retained, the deformation and summation
are legitimate term by term.
\end{proof}

When the level-$m$ inherited points are simple square-root germs, their
combined contribution has the particularly transparent form
\begin{equation}\label{eq:root-level-transseries}
 \rho^{-n/m}n^{-3/2}
 \sum_{j=0}^{m-1}e^{-2\pi\ii jn/m}
 \sum_{r\ge0}\frac{A_{m,j,r}}{n^r}.
\end{equation}
Relative to the dominant sector, its exponential smallness is
\begin{equation}\label{eq:relative-root-scale}
 \exp\!\left[-\lambda\left(1-\frac1m\right)n\right].
\end{equation}
The distinction between the Poincar\'e sector and exponentially small
sectors follows the transseries conventions emphasized in
\cite{ProveItTransseries}.  The first inherited level, $m=2$, is therefore of relative order
$e^{-\lambda n/2}$, beyond every inverse power of $n$.

\subsection{Recursive construction of every local germ}

The singularity recursion can be implemented entirely with the same Bell
polynomials used earlier.  Let $z=\sigma(1-u)$.  If $\tau=\sigma^k$ is a
parent singularity, then
\begin{equation}\label{eq:Pk-def}
 1-\frac{z^k}{\tau}=1-(1-u)^k
 =ku\left(1+\sum_{r\ge1}p_{k,r}u^r\right),
 \qquad
 p_{k,r}=\frac{(-1)^r}{k}\binom{k}{r+1}.
\end{equation}
For a parent monomial $(1-z^k/\tau)^\beta$,
\begin{align}
 \left(1-\frac{z^k}{\tau}\right)^\beta
 ={}&(ku)^\beta
 \sum_{N\ge0}u^N
 \sum_{j=0}^{N}\binom\beta j
 \Bhat_{N,j}(p_{k,1},p_{k,2},\ldots).
 \label{eq:parent-propagation}
\end{align}
Thus every inherited local coefficient is a finite ordinary-Bell sum of
parent coefficients.

After all singular summands $T(z^k)/k$ have been propagated, add the regular
Taylor parts to obtain the local generalized series for $R(z)$.  Exponentiate
it to form $\zeta(z)=z e^{R(z)}$; after replacing $u$ by a common ramified
coordinate $w=u^{1/L}$, this exponentiation is exactly
\eqref{eq:complete-bell-exp}.

There are then two outer-composition cases.

\paragraph{Noncritical inherited point.}
If $\zeta(\sigma)\ne e^{-1}$ and the selected Cayley branch is regular there,
write $\Delta(u)=\zeta(z)-\zeta(\sigma)$.  Then
\begin{equation}\label{eq:noncritical-C-composition}
 T(z)=\sum_{r\ge0}\frac{\Cay^{(r)}(\zeta(\sigma))}{r!}\Delta(u)^r,
\end{equation}
and powers of $\Delta$ are computed by ordinary Bell polynomials.  The
Cayley derivatives themselves may be generated from
\begin{equation}\label{eq:Cay-derivative}
 \Cay'(w)=\frac{\Cay(w)}{w(1-\Cay(w))}
\end{equation}
by differentiation and Fa\`a di Bruno calculus.

\paragraph{Critical point.}
If $\zeta(\sigma)=e^{-1}$, define $q_\sigma(u)=1-e\zeta(z)$.  Compose its
leading ramified series with the universal Cayley germ
\begin{equation}\label{eq:critical-C-composition}
 T(z)=1+\sum_{r\ge1}c_rq_\sigma(u)^{r/2}.
\end{equation}
The coefficient extraction is the ramified version of
\eqref{eq:t-n-bell}.  If the leading exponent of $q_\sigma$ is $\delta$,
then a generic critical composition begins with exponent $\delta/2$.
Hence a coincidence between an inherited square root and a Cayley critical
point can create quarter powers; \cref{thm:general-sector} already allows
this.

\begin{algorithm}[Local-to-global transseries engine]\label{alg:global-engine}
To compute every sector inside a prescribed radius $R<1$ and to a prescribed
power order:
\begin{enumerate}
\item Enumerate accessible inherited preimages and continued critical
      solutions in $|z|<R$, recording the parent relation and branch.
\item At each point $\sigma$, choose a common ramification index $L$ and
      propagate all parent germs by \eqref{eq:parent-propagation}.
\item Sum the propagated germs in $R(z)$ and exponentiate with complete Bell
      polynomials to obtain the local germ of $\zeta$.
\item Use \eqref{eq:noncritical-C-composition} or
      \eqref{eq:critical-C-composition} to obtain the local germ of $T$.
\item Transfer every noninteger monomial with
      \eqref{eq:general-sector-coeff}; pair conjugate sectors to obtain a
      manifestly real expression on integer $n$.
\end{enumerate}
Every operation at a fixed truncation order is finite and triangular.
\end{algorithm}

\begin{frontierbox}
The all-orders dominant sector is completely explicit and unconditional.
The recursive engine in \cref{alg:global-engine} is a constructive
beyond-all-orders completion: once a finite set of continued singular germs
has been certified, every coefficient in their exponential sectors follows
from finite Bell-polynomial operations.  A natural next theorem would prove
$\mathsf H(R)$ uniformly as $R\uparrow1$ and classify all critical
coincidences on the resulting Riemann surface.
\end{frontierbox}

\section{Beyond-all-orders inversion}\label{sec:inverse-exponential}

\subsection{A chosen transseries interpolation}

For noninteger $x$, an oscillatory factor $\sigma^{-x}$ requires a branch of
$\Log\sigma$.  Different branches agree at integer $x=n$ but generally
differ between integers.  Therefore the exponentially small part of a
\emph{functional} inverse is not canonical until a continuation convention
has been chosen.

Fix such a convention and pair conjugate branches so that the interpolation
is real on the positive real axis.  Write the completed forward transseries
as
\begin{equation}\label{eq:full-interpolation}
 \mathfrak a(x)=\mathcal A_0(x)\bigl(1+\mathcal E(x)\bigr),
 \qquad
 \mathcal A_0(x)=Ce^{\lambda x}x^{-p}S(x^{-1}),
\end{equation}
where $\mathcal E$ is a sum of exponentially small sectors.  A term from a
singularity $\sigma=|\sigma|e^{\ii\theta}$ and local exponent $\beta$ has,
relative to $\mathcal A_0$, the scale
\begin{equation}\label{eq:relative-general-sector}
 Q_{\sigma,\beta}(x)
 :=\exp\!\bigl[-\mu_\sigma x-\ii\theta x\bigr]
 x^{p-\beta-1},
 \qquad
 \mu_\sigma:=\log\frac{|\sigma|}{\rho}>0,
\end{equation}
multiplied by a descending power series.  For
$\sigma=\rho^{1/m}e^{2\pi\ii j/m}$ and a square-root exponent
$\beta=1/2$, this reduces to
\begin{equation}\label{eq:root-inverse-monomial}
 Q_{m,j}(x)=
 e^{-\lambda(1-1/m)x}e^{-2\pi\ii jx/m},
\end{equation}
with no extra leading power of $x$.

\subsection{Carrier-first implicit equation}

Let $\nu_0(y)$ be the complete algebraic inverse from
\cref{sec:inverse-carrier}; thus $\mathcal A_0(\nu_0)=y$.  Put
\begin{equation}\label{eq:delta-exp-def}
 \nu_{\mathrm{full}}(y)=\nu_0(y)+\delta(y).
\end{equation}
With
\begin{equation}\label{eq:F0-def}
 F_0(x):=\log\mathcal A_0(x)
 =\log C+\lambda x-p\log x+H(x^{-1}),
\end{equation}
the exact transseries correction equation is
\begin{equation}\label{eq:delta-exact}
 \boxed{
 F_0(\nu_0+\delta)-F_0(\nu_0)
 +\log\bigl(1+\mathcal E(\nu_0+\delta)\bigr)=0.}
\end{equation}
The derivative
\begin{equation}\label{eq:F0-prime}
 F_0'(x)=\lambda-\frac px-\frac{H'(x^{-1})}{x^2}
 =\lambda+\Ord(x^{-1})
\end{equation}
is invertible in the slow coefficient field.

\begin{theorem}[General coefficient recursion for the inverse transseries]
\label{thm:inverse-exponential-recursion}
Choose finitely many independent small transmonomials
$\bm Q=(Q_1,\ldots,Q_s)$ and truncate the forward transseries accordingly.
There is a unique formal solution
\begin{equation}\label{eq:delta-multi}
 \delta=\sum_{\bm n\in\Nzero^s\setminus\{\bm0\}}
 D_{\bm n}(\nu_0^{-1},\log\nu_0)\bm Q^{\bm n}
\end{equation}
of \eqref{eq:delta-exact}, ordered by total transmonomial valuation.  Its
coefficients obey the triangular recursion
\begin{equation}\label{eq:Dmulti-recursion}
 \boxed{
 D_{\bm n}=-\frac1{F_0'(\nu_0)}
 [\bm Q^{\bm n}]
 \left\{
 \sum_{r\ge2}\frac{F_0^{(r)}(\nu_0)}{r!}\delta^r
 +\log\bigl(1+\mathcal E(\nu_0+\delta)\bigr)
 \right\},}
\end{equation}
where the coefficient on the right is evaluated using only terms of
$\delta$ of strictly smaller transmonomial valuation.  Every
$D_{\bm n}$ then has a computable descending power-logarithmic expansion in
$\nu_0$.
\end{theorem}

\begin{proof}
Expand the first term of \eqref{eq:delta-exact} at $\nu_0$:
\[
 F_0(\nu_0+\delta)-F_0(\nu_0)
 =F_0'(\nu_0)\delta+
  \sum_{r\ge2}\frac{F_0^{(r)}(\nu_0)}{r!}\delta^r.
\]
Since $F_0'(\nu_0)$ has nonzero leading term $\lambda$, it is invertible.
At any nonzero multiindex, the coefficient of the linear term is
$F_0'(\nu_0)D_{\bm n}$.  Every other occurrence of $\delta$ is either
nonlinear or multiplied by a positive transmonomial from $\mathcal E$, so it
contains only lower coefficients in the well-ordered support.  Solving the
linear equation gives \eqref{eq:Dmulti-recursion}.  The slow-field expansion
uses the same Bell and composition recursions as \cref{sec:inverse-pl}.
\end{proof}

The first exponential correction is especially simple.  If
$\mathcal E^{[1]}$ denotes the part linear in the chosen small monomials,
then
\begin{equation}\label{eq:first-exp-inverse}
 \boxed{
 \delta^{[1]}(y)=
 -\frac{\mathcal E^{[1]}(\nu_0(y))}{F_0'(\nu_0(y))}.}
\end{equation}
Higher sectors include products of forward sectors, derivatives caused by
the shift $x\mapsto x+\delta$, and the logarithmic cumulants of
$1+\mathcal E$; \eqref{eq:Dmulti-recursion} accounts for all of them without
case-by-case guessing.

\subsection{Canonical and noncanonical information}

The inverse hierarchy has three layers.
\begin{enumerate}
\item The Lambert carrier and its algebraic corrections are canonical and
      depend only on the dominant Poincar\'e sector.
\item Exponentially small smooth corrections are canonical only after the
      analytic continuation and logarithm branches defining the
      interpolation have been fixed.
\item The exact discrete generalized inverse additionally requires the
      lattice rounding described in \cref{sec:discrete}.
\end{enumerate}
This hierarchy is not a defect.  It records exactly how much information the
integer sequence itself determines at each asymptotic scale.

\section{Coefficient algorithms and reproducibility}\label{sec:algorithms}

\subsection{Exact sequence generation}

The following sieve form of \eqref{eq:exact-recurrence} avoids recomputing
divisors.

\begin{lstlisting}[caption={Exact generation of $a_1,\ldots,a_N$ (language-neutral pseudocode).},label={lst:exact-sequence}]
a[0] := 0
b[0..N] := 0
a[1] := 1
for multiple = 1..N:
    b[multiple] += 1              # contribution 1*a[1]

for n = 2..N:
    numerator := sum_{k=1}^{n-1} b[k] * a[n-k]
    assert numerator mod (n-1) = 0
    a[n] := numerator / (n-1)
    for multiple = n, 2n, 3n, ... <= N:
        b[multiple] += n * a[n]
\end{lstlisting}
The integer divisibility assertion is a useful exact check on every newly
computed term.

\subsection{Solving for $\rho$ without evaluating $T$ at its boundary}

Equation \eqref{eq:rho-scalar} uses only $T(r^k)$ with $k\ge2$.  Even when
$r$ is extremely close to $\rho$, these arguments remain uniformly
subcritical.  A certified procedure is therefore:
\begin{enumerate}
\item generate enough exact $a_n$ by \cref{lst:exact-sequence};
\item evaluate each $T(r^k)$ by a truncated positive series with a tail
      majorant;
\item use interval Newton or bisection on
      $f(r)=\log r+1+\sum_{k\ge2}T(r^k)/k$;
\item verify $f'(r)>0$ on the final interval.
\end{enumerate}
This is much more stable than attempting to estimate the singularity from
ratios $a_{n+1}/a_n$, whose convergence carries inverse-power corrections.

\subsection{Power, exponential, and logarithm recurrences}

Bell polynomials are conceptually clean, but dense recurrences are faster in
software.  If
\[
 F(u)=1+\sum_{n\ge1}f_nu^n,
 \qquad F(u)^\gamma=\sum_{n\ge0}v_nu^n,
\]
then differentiation of $Fv'=\gamma vF'$ gives
\begin{equation}\label{eq:power-series-recurrence}
 v_0=1,\qquad
 v_n=\frac1n\sum_{k=1}^{n}\bigl((\gamma+1)k-n\bigr)f_kv_{n-k}.
\end{equation}
If $A(u)=\sum_{n\ge0}A_nu^n$ and
$E(u)=e^{A(u)}=\sum E_nu^n$, then
\begin{equation}\label{eq:exp-series-recurrence}
 E_0=e^{A_0},\qquad
 E_n=\frac1n\sum_{k=1}^{n}kA_kE_{n-k}.
\end{equation}
If $F(0)=1$ and $H=\log F$, then
\begin{equation}\label{eq:log-series-recurrence-general}
 H_n=f_n-\frac1n\sum_{k=1}^{n-1}kH_kf_{n-k}.
\end{equation}
These recurrences are exactly equivalent to the Bell-polynomial formulas and
are convenient for high-order numerical work.

\subsection{End-to-end dominant-sector pipeline}

\begin{algorithm}[All coefficients through order $M$]\label{alg:dominant-pipeline}
Given a target order $M$:
\begin{enumerate}
\item Generate exact tree numbers and solve \eqref{eq:rho-scalar} for $\rho$.
\item Compute the coefficients of $R(\rho(1-s))$ through $s^{2M+1}$ using
      \eqref{eq:R-direct-coeff}; exponentiate by
      \eqref{eq:exp-series-recurrence} to obtain $q_1,\ldots,q_{2M+1}$.
\item Compute $c_1,\ldots,c_{2M+1}$ from
      \eqref{eq:c-r-lagrange}, using
      \eqref{eq:power-series-recurrence} for $\Phi^{-r/2}$.
\item Compute $t_1,\ldots,t_{2M+1}$ by \eqref{eq:t-n-coeff}.
\item Compute $g_r(j+1/2)$ from \eqref{eq:g-recurrence} and form
      $\tau_0,\ldots,\tau_M$ by \eqref{eq:tau-master}.
\item Normalize to $s_m=\tau_m/\tau_0$, compute $h_m$ by
      \eqref{eq:h-recurrence}, then compute $d_m$ by
      \eqref{eq:d-coeff-extraction} or $P_m$ by
      \eqref{eq:P-coeff-recurrence}.
\end{enumerate}
At every stage, the coefficient of order $m$ depends only on data of order
at most $m$; no nonlinear global solve is required after $\rho$ is known.
\end{algorithm}

\subsection{Symbolic and numerical audit checks}

A robust implementation should verify all of the following independent
identities.
\begin{enumerate}
\item The exact recurrence reproduces the known initial values
      $1,1,2,4,9,20,\ldots$.
\item The critical residuals $T(\rho)-1$ and $e\zeta(\rho)-1$ vanish within
      the certified numerical enclosure.
\item The identity
      $1-e\zeta(\rho(1-s))=q(s)$ holds coefficientwise to the requested
      order.
\item Substituting the computed $t_n$ into
      $T=\zeta e^T$ leaves a Puiseux residual beyond the requested order.
\item The two formulas \eqref{eq:tau-master} and direct expansion of the
      exact transfer coefficient \eqref{eq:exact-transfer} agree.
\item Substitution of the $d_m$ into \eqref{eq:D-equation}, and of the
      $P_m$ into \eqref{eq:R-PL-equation}, leaves zero through the requested
      order.
\end{enumerate}
These checks separate derivation errors from loss of numerical precision.

\section{Numerical validation}\label{sec:numerics}

\subsection{Forward expansion}

Let
\begin{equation}\label{eq:A-M}
 A_M(n):=\frac{\rho^{-n}}{\sqrt\pi\,n^{3/2}}
 \sum_{m=0}^{M-1}\frac{\tau_m}{n^m}.
\end{equation}
\Cref{tab:forward-errors} compares $A_M(n)$ with exact integers generated by
\eqref{eq:exact-recurrence}.  The entries are relative errors
$|A_M(n)/a_n-1|$.

\begin{table}[H]
\centering
\caption{Relative errors of the dominant asymptotic expansion.}
\label{tab:forward-errors}
\begin{tabular}{@{}rccccc@{}}
\toprule
$n$ & $M=1$ & $M=4$ & $M=8$ & $M=12$ \\
\midrule
20  & $6.5807\times10^{-3}$ & $9.7082\times10^{-5}$ & $6.3781\times10^{-6}$ & $1.5969\times10^{-5}$\\
50  & $2.2224\times10^{-3}$ & $1.9668\times10^{-6}$ & $3.1903\times10^{-9}$ & $5.8824\times10^{-11}$\\
100 & $1.0554\times10^{-3}$ & $1.1313\times10^{-7}$ & $1.0129\times10^{-11}$ & $9.7958\times10^{-15}$\\
200 & $5.1460\times10^{-4}$ & $6.8101\times10^{-9}$ & $3.6381\times10^{-14}$ & $2.0834\times10^{-18}$\\
500 & $2.0280\times10^{-4}$ & $1.7064\times10^{-10}$ & $2.2767\times10^{-17}$ & $3.2504\times10^{-23}$\\
\bottomrule
\end{tabular}
\end{table}

At $n=20$, twelve terms already pass the least term and perform worse than
eight terms.  At larger $n$, the optimal truncation moves outward exactly as
expected for a divergent asymptotic expansion.

\subsection{Inverse expansion at exact sequence values}

For $y=a_n$, define
\begin{equation}\label{eq:nu-K}
 \nu_K(y):=x_0(y)+\sum_{m=1}^{K}\frac{d_m}{x_0(y)^m}.
\end{equation}
\Cref{tab:inverse-errors} records the absolute index error
$|\nu_K(a_n)-n|$.  The ``carrier'' column means $K=0$.

\begin{table}[H]
\centering
\caption{Absolute errors of the Lambert-carrier inverse.}
\label{tab:inverse-errors}
\begin{tabular}{@{}rccccc@{}}
\toprule
$n$ & carrier & $K=1$ & $K=4$ & $K=8$ & $K=12$\\
\midrule
20  & $6.5450\times10^{-3}$ & $1.9143\times10^{-3}$ & $3.7564\times10^{-5}$ & $9.3870\times10^{-6}$ & $1.7523\times10^{-5}$\\
50  & $2.1114\times10^{-3}$ & $2.5861\times10^{-4}$ & $3.3169\times10^{-7}$ & $9.9830\times10^{-10}$ & $2.7265\times10^{-11}$\\
100 & $9.8802\times10^{-4}$ & $6.1593\times10^{-5}$ & $9.3442\times10^{-9}$ & $1.5282\times10^{-12}$ & $2.1783\times10^{-15}$\\
200 & $4.7827\times10^{-4}$ & $1.5047\times10^{-5}$ & $2.7890\times10^{-10}$ & $2.7117\times10^{-15}$ & $2.2840\times10^{-19}$\\
500 & $1.8766\times10^{-4}$ & $2.3752\times10^{-6}$ & $2.7824\times10^{-12}$ & $6.7465\times10^{-19}$ & $1.4159\times10^{-24}$\\
\bottomrule
\end{tabular}
\end{table}

The inverse error is approximately the logarithmic forward residual divided
by $\lambda$, in agreement with \eqref{eq:inverse-error}.  The same
least-term behavior is visible at $n=20$.

\subsection{Independent recovery of the Puiseux coefficients}

A direct implementation of \eqref{eq:R-direct-coeff} with exact $a_n$,
followed by \eqref{eq:exp-series-recurrence},
\eqref{eq:c-r-lagrange}, and \eqref{eq:t-n-coeff}, gives
\[
 t_1=-1.55949002037464088554220739589\ldots,
\]
\[
 t_9=-0.1952123835975564636127699\ldots,
 \qquad
 t_{17}=-0.1445471935709097054769461\ldots.
\]
Substitution into \eqref{eq:tau-master} gives, independently,
\[
 \tau_0=0.7797450101873204427711037\ldots,
 \qquad
 \tau_8=67342.06920114653048340563\ldots.
\]
These values agree with the published coefficient tables
\cite{Genitrini2016} to the accuracy justified by their input precision;
the displayed values here were recomputed with higher internal precision.

\section{Further deductions and research directions}\label{sec:research}

\subsection{Large-order behavior of $\tau_m$}

The data strongly suggest factorial-over-power growth of the algebraic
coefficients.  A precise late-term law should be controlled by the nearest
singular obstruction of the local $s^{1/2}$ expansion and, after coefficient
transfer, by the first adjacent exponential sectors.  The singularity engine
of \cref{alg:global-engine} gives a concrete route to this law:
compute the first continued germs, derive their Stokes constants relative to
the dominant germ, and match the large-$m$ growth of $\tau_m$.  This would
turn the present finite-level transseries into a quantitative resurgence
statement.

\subsection{A root-sector conjecture}

\begin{conjecture}[Generic inherited half-power skeleton]
On noncritical continued sheets, every point
$\rho^{1/m}e^{2\pi\ii j/m}$ that is reached without an earlier obstruction
carries a square-root germ inherited from $z^m=\rho$.  Its coefficient
sector therefore has the form \eqref{eq:root-level-transseries}.  Critical
coincidences are isolated and are the only source of finer ramification on
the inherited skeleton.
\end{conjecture}

This conjecture is local and testable: interval analytic continuation along
a prescribed path can certify both accessibility and the noncritical
condition $\zeta(\sigma)\ne e^{-1}$.

\subsection{Stokes data and lateral continuation}

The positive singularity $\rho$ lies on the natural direction of coefficient
extraction.  To expose sectors farther along related cuts one should compare
upper and lower lateral continuations.  Their difference is a discontinuity
transseries whose coefficients can be propagated by the same local Bell
calculus.  Pairing lateral values is also the natural way to obtain real
median interpolations for the inverse problem.

\subsection{Formalization targets}

The development separates into layers suitable for mechanization.
\begin{enumerate}
\item Purely algebraic identities: ordinary and exponential Bell
      polynomials, the Cayley coefficient formula, gamma-ratio recurrences,
      and the triangular inverse recursions.
\item Local analytic theorems: implicit square-root inversion, convergence
      of the subcritical derivative sums, and fixed-order transfer.
\item Global analytic certificates: continuation domains, singularity
      enumeration, and contour deformation for $\mathsf H(R)$.
\end{enumerate}
The first layer is especially well suited to exact rational and algebraic
computation, while the second can use interval bounds around numerical
constants such as $\rho$.

\section{Conclusion}

The Butcher--P\'olya rooted-tree numbers admit a particularly clean
all-orders calculus once the generating function is factored through the
Cayley tree function.  The complete chain is
\[
 \text{subcritical values of }T
 \longrightarrow \zeta^{(m)}(\rho)
 \longrightarrow q_m
 \longrightarrow t_m
 \longrightarrow \tau_m
 \longrightarrow h_m
 \longrightarrow d_m\text{ or }P_m.
\]
Every arrow is given by a finite Bell-polynomial or Bernoulli-polynomial
formula.  The dominant result is the compact convolution
\eqref{eq:tau-master}; the functional inverse is the exact Lambert carrier
\eqref{eq:carrier-W-tree} with the all-orders correction
\eqref{eq:d-coeff-extraction}, or equivalently the polynomial--logarithmic
recursion \eqref{eq:P-coeff-recurrence}.

Beyond all powers, the recursive generating equation supplies a singularity
grammar.  The universal sector formula \eqref{eq:general-sector-coeff} and
the local propagation algorithm turn any certified finite continuation into
a coefficient transseries, while \eqref{eq:Dmulti-recursion} reverses that
transseries around the fully normalized carrier.  The remaining global
frontier is no longer a question of missing coefficient algebra; it is the
analytic classification of accessible sheets, critical preimages, and their
Stokes connections.

\appendix

\section{Analytic details at the dominant singularity}\label{app:analytic}

This appendix records the local analytic argument in a form that makes clear
which properties are specific to P\'olya trees and which are universal.

\begin{proposition}[Analyticity of the perturbation]\label{prop:R-analytic}
If $\rho$ is the radius of convergence of $T$, then
\[
 R(z)=\sum_{k\ge2}\frac{T(z^k)}k
\]
is analytic in $|z|<\sqrt\rho$.  On every compact subdisc, the series and all
of its derivatives converge normally.
\end{proposition}

\begin{proof}
Fix $r<\sqrt\rho$.  Then $r^2<\rho$.  Since $T(w)=w+\Ord(w^2)$ at the
origin, there are constants $M$ and $r_0<1$ such that
$|T(w)|\le M|w|$ for $|w|\le r_0$.  Split the sum into finitely many $k$
for which $r^k>r_0$ and the remaining tail.  In the tail,
\[
 \left|\frac{T(z^k)}k\right|\le\frac{M r^k}{k},
\]
which is summable.  The same argument after differentiating uses a polynomial
factor in $k$ multiplied by $r^{k-m}$, still summable.  Weierstrass normal
convergence proves analyticity and termwise differentiation.
\end{proof}

\begin{proposition}[Critical endpoint and transversality]
\label{prop:critical-transverse}
Along the positive axis, $\zeta(r)$ is strictly increasing.  At the first
singularity $r=\rho$,
\[
 \zeta(\rho)=e^{-1},\qquad T(\rho)=1,
 \qquad \zeta'(\rho)>0.
\]
Consequently $q_1=e\rho\zeta'(\rho)>0$.
\end{proposition}

\begin{proof}
The coefficients of $T$, hence of $R$ and $\zeta=z e^R$, are nonnegative;
the coefficient of $z$ in $\zeta$ is positive.  Therefore $\zeta'(r)>0$
for $r>0$ in its interval of analyticity.  While $\zeta(r)<e^{-1}$, the
principal Cayley function $\Cay(\zeta(r))$ is analytic and solves the
functional equation.  If $\zeta$ did not reach $e^{-1}$ at $\rho$, the
composition $T=\Cay\circ\zeta$ would extend analytically beyond $\rho$,
contradicting the definition of the radius.  Thus
$\zeta(\rho)=e^{-1}$, and \eqref{eq:cayley-def} gives $T(\rho)=1$.
Strict positivity of $\zeta'$ gives the last assertion.
\end{proof}

\begin{proposition}[Uniqueness on the dominant circle]
\label{prop:unique-dominant}
The positive point $z=\rho$ is the only singularity of the branch defined at
zero on $|z|=\rho$.
\end{proposition}

\begin{proof}
The series $\zeta(z)$ has nonnegative coefficients.  Its first relevant
terms are
\[
 \zeta(z)=z+\frac12z^3+\frac13z^4+\cdots,
\]
because $R(z)=z^2/2+z^3/3+\cdots$.  For $|z|=\rho$, the triangle inequality
gives $|\zeta(z)|\le\zeta(\rho)=e^{-1}$.  Equality would require the nonzero
terms $z,z^3,z^4$ to have the same argument.  Writing $z=\rho e^{\ii\theta}$,
this implies $e^{2\ii\theta}=e^{3\ii\theta}=1$, hence
$e^{\ii\theta}=1$ and $z=\rho$.  At every other point the inequality is
strict, so the principal Cayley function is analytic at $\zeta(z)$.  The
perturbation $\zeta$ is analytic on the larger disc $|z|<\sqrt\rho$ by
\cref{prop:R-analytic}; therefore $T=\Cay\circ\zeta$ is analytic at every
other point of $|z|=\rho$.
\end{proof}

\begin{theorem}[Local square-root expansion with analytic remainder]
\label{thm:local-square-root-detailed}
There is a dented neighborhood of $\rho$ in which
\[
 T(z)=1+\sum_{n=1}^{N}t_n\left(1-\frac z\rho\right)^{n/2}
 +\Ord\!\left(\left|1-\frac z\rho\right|^{(N+1)/2}\right)
\]
for every $N\ge1$, with coefficients given by
\cref{thm:t-n-formula}.
\end{theorem}

\begin{proof}
By \cref{prop:R-analytic}, $q(s)$ in \eqref{eq:q-def} is analytic near
$s=0$; by \cref{prop:critical-transverse}, it has a simple zero with positive
coefficient $q_1$.  Therefore $q(s)=q_1s(1+U(s))$, where $U$ is analytic and
$U(0)=0$.  Choose a slit neighborhood of $s=0$ and a consistent square root.
The convergent Cayley Puiseux series \eqref{eq:cayley-puiseux} may be
composed with this analytic map.  Ordinary convergence of the composition
and Taylor's theorem give the displayed remainder.  Coefficient extraction
in the convergent composition is exactly \eqref{eq:t-n-bell}.
\end{proof}

Together, \cref{prop:unique-dominant,thm:local-square-root-detailed} supply
the standard Delta-domain hypotheses used in the proof of
\cref{thm:dominant-expansion}.

\section{Coefficient tables for inversion}\label{app:inverse-tables}

\begingroup
\small
\begin{longtable}{@{}r@{\enspace}r@{\enspace}r@{\enspace}r@{}}
\caption{Normalized forward, logarithmic, and Lambert-correction
coefficients.}
\label{tab:inverse-coefficients}\\
\toprule
$m$ & $s_m$ & $h_m$ & $d_m$\\
\midrule
\endfirsthead
\toprule
$m$ & $s_m$ & $h_m$ & $d_m$\\
\midrule
\endhead
1 & $0.1004034800964014348$ & $0.1004034800964014348$ & $-0.09264385352561312984$\\
2 & $0.5039343151023035331$ & $0.4988938856945692936$ & $-0.5885630399313472628$\\
3 & $1.972284908382278469$ & $1.922025533841821822$ & $-2.596680167407531021$\\
4 & $10.51734090413157399$ & $10.19739642337376887$ & $-13.14905339996706569$\\
5 & $73.47647498402047106$ & $71.47146706207073604$ & $-85.49870151462556652$\\
6 & $645.1397552456612928$ & $630.8599376073372479$ & $-711.2629148329465609$\\
7 & $6873.530274463408620$ & $6753.629059387619229$ & $-7306.788270284967257$\\
8 & $86364.21948371140798$ & $85171.04383262353790$ & $-89561.76088263724860$\\
9 & $1.250954458605978292\times10^6$ & $1.236987135642947527\times10^6$ & $-1.274829503083708903\times10^6$\\
10 & $2.051559456480361859\times10^7$ & $2.032577710708850382\times10^7$ & $-2.064109900032138585\times10^7$\\
\bottomrule
\end{longtable}
\endgroup

The rapid growth does not contradict asymptotic validity.  The series are to
be truncated according to the large parameter, and the usable order grows
with $n$ or $x_0$.

\section{Derivation of selected inverse coefficients}\label{app:d-derivations}

For illustration, expand \eqref{eq:D-equation} through $t^4$.  Since
$D=d_1t+d_2t^2+d_3t^3+d_4t^4+\cdots$, one has
\begin{align*}
 \log(1+tD)
 &=d_1t^2+d_2t^3+
   \left(d_3-\frac12d_1^2\right)t^4+\Ord(t^5),\\
 \frac{t}{1+tD}
 &=t-d_1t^3-d_2t^4+\Ord(t^5),\\
 H\!\left(\frac{t}{1+tD}\right)
 &=h_1t+h_2t^2+(h_3-h_1d_1)t^3\\
 &\quad +(h_4-2h_2d_1-h_1d_2)t^4+\Ord(t^5).
\end{align*}
Equating powers of $t$ gives
\begin{align*}
 \lambda d_1+h_1&=0,\\
 \lambda d_2-pd_1+h_2&=0,\\
 \lambda d_3-pd_2+h_3-h_1d_1&=0,\\
 \lambda d_4-p\left(d_3-\frac12d_1^2\right)
 +h_4-2h_2d_1-h_1d_2&=0.
\end{align*}
Successive substitution yields \eqref{eq:d1}--\eqref{eq:d4}.  This small
calculation displays the triangular mechanism that the general Bell formula
\eqref{eq:d-bell-first}--\eqref{eq:d-bell-second} packages at arbitrary
order.

For completeness, the next coefficient is
\begin{align}
 d_5=-\frac1{2\lambda^5}\bigl(&4\lambda^2h_1^3
 +9\lambda p^2h_1^2+14\lambda^2p h_1h_2
 +8\lambda^3h_1h_3+2p^4h_1\notag\\
 &+4\lambda^3h_2^2+2\lambda p^3h_2
 +2\lambda^2p^2h_3+2\lambda^3p h_4
 +2\lambda^4h_5\bigr).
\label{eq:d5}
\end{align}

\section{A compact formula sheet}\label{app:formula-sheet}

For reference, the complete dominant and inverse pipeline is collected here.
\begin{align*}
 T(z)&=z\exp\!\left(\sum_{k\ge1}\frac{T(z^k)}k\right),\\
 \log\rho+1+\sum_{k\ge2}\frac{T(\rho^k)}k&=0,\\
 q_m&=\frac{(-1)^{m+1}e\rho^m}{m!}\zeta^{(m)}(\rho),\\
 c_r&=-\frac{2^{r/2}}r[v^{r-1}]
 \left(\frac{2(1-(1-v)e^v)}{v^2}\right)^{-r/2},\\
 t_n&=\sum_{\substack{1\le r\le n\\r\equiv n\ (2)}}
 c_rq_1^{r/2}[s^{(n-r)/2}](1+U(s))^{r/2},\\
 \kappa_m(\beta)&=\frac{(-1)^{m+1}}{m(m+1)}
 \bigl(B_{m+1}(-\beta)-B_{m+1}(1)\bigr),\\
 g_r(\beta)&=\frac1{r!}Y_r(1!\kappa_1,\ldots,r!\kappa_r),\\
 \tau_m&=\sum_{j=0}^{m}
 \frac{(-1)^{j+1}(2j+1)!!}{2^{j+1}}
 t_{2j+1}g_{m-j}(j+\tfrac12),\\
 h_m&=\sum_{k=1}^{m}\frac{(-1)^{k-1}}k
 \Bhat_{m,k}(s_1,s_2,\ldots),\\
 x_0(y)&=-\frac{3}{2\lambda}W_{-1}\!\left(
 -\frac{2\lambda}{3}(C/y)^{2/3}\right),\\
 d_m&=[t^m]\left\{\frac p\lambda\log(1+tD_{<m})
 -\frac1\lambda H\!\left(\frac{t}{1+tD_{<m}}\right)\right\},\\
 \nu(y)&\sim x_0(y)+\sum_{m\ge1}d_mx_0(y)^{-m}.
\end{align*}
For a general singular germ $c_{\sigma,\beta}(1-z/\sigma)^\beta$,
replace the dominant formula by the universal sector coefficient
\[
 A_{\sigma,\beta,r}=
 \frac{c_{\sigma,\beta}}{\Gamma(-\beta)}g_r(\beta).
\]

\begin{thebibliography}{99}

\bibitem{Butcher1972}
J.~C. Butcher,
\newblock An algebraic theory of integration methods,
\newblock \emph{Mathematics of Computation} \textbf{26} (1972), 79--106.

\bibitem{Comtet1974}
L.~Comtet,
\newblock \emph{Advanced Combinatorics: The Art of Finite and Infinite Expansions},
\newblock D. Reidel, Dordrecht, 1974.

\bibitem{Corless1996}
R.~M. Corless, G.~H. Gonnet, D.~E.~G. Hare, D.~J. Jeffrey, and D.~E. Knuth,
\newblock On the Lambert $W$ function,
\newblock \emph{Advances in Computational Mathematics} \textbf{5} (1996),
329--359.

\bibitem{DLMF}
NIST Digital Library of Mathematical Functions,
\newblock Chapters 5 and 24: gamma functions, Bernoulli polynomials, and
asymptotic expansions,
\newblock \url{https://dlmf.nist.gov/}, accessed September 2026.

\bibitem{FlajoletSedgewick2009}
P.~Flajolet and R.~Sedgewick,
\newblock \emph{Analytic Combinatorics},
\newblock Cambridge University Press, 2009.

\bibitem{Genitrini2016}
A.~Genitrini,
\newblock Full asymptotic expansion for P\'olya structures,
\newblock arXiv:1605.00837 [math.CO], 2016,
\newblock \url{https://arxiv.org/abs/1605.00837}.

\bibitem{McLachlan2017}
R.~I. McLachlan, K.~Modin, H.~Munthe-Kaas, and O.~Verdier,
\newblock What are Butcher series, really? The story of rooted trees and
numerical methods for evolution equations,
\newblock \emph{BIT Numerical Mathematics} \textbf{57} (2017), 477--501.

\bibitem{OEISA000081}
OEIS Foundation Inc.,
\newblock A000081: number of unlabeled rooted trees with $n$ nodes,
\newblock \url{https://oeis.org/A000081}, accessed September 2026.

\bibitem{OEISOtterConstants}
OEIS Foundation Inc.,
\newblock A051491 and A187770: Otter's rooted-tree growth and asymptotic
constants,
\newblock \url{https://oeis.org/A051491},
\url{https://oeis.org/A187770}, accessed September 2026.

\bibitem{Otter1948}
R.~Otter,
\newblock The number of trees,
\newblock \emph{Annals of Mathematics} \textbf{49} (1948), 583--599.

\bibitem{Polya1937}
G.~P\'olya,
\newblock Kombinatorische Anzahlbestimmungen f\"ur Gruppen, Graphen und
chemische Verbindungen,
\newblock \emph{Acta Mathematica} \textbf{68} (1937), 145--254.

\bibitem{ProveItCoefficientCalculus}
V.~Reshetnikov, ProveIt project,
\newblock \emph{Combinatorial Coefficient Calculus: Stirling, Bell,
Eulerian, and Bernoulli--Euler Families; Bell-Polynomial and Fa\`a di Bruno
Calculus; Lagrange--B\"urmann and Multivariate Good Inversion},
\newblock research draft, September 2026,
\newblock \url{https://github.com/VladimirReshetnikov/ProveIt/tree/main/Analysis/FabiusFunction/docs/semi-formalized-research-frontiers/drafts/combinatorial-coefficient-calculus/Combinatorial_Coefficient_Calculus}.

\bibitem{ProveItTransseries}
V.~Reshetnikov, ProveIt project,
\newblock Series-and-transseries research drafts, including
\emph{Polynomial--Logarithmic Transseries} and special-function inversion,
\newblock September 2026,
\newblock \url{https://github.com/VladimirReshetnikov/ProveIt/tree/main/Analysis/FabiusFunction/docs/semi-formalized-research-frontiers/drafts/series-and-transseries}.

\end{thebibliography}

\end{document}
